% Created 2026-05-18 Mon 08:51
% Intended LaTeX compiler: pdflatex
\documentclass[11pt]{article}

\usepackage[utf8]{inputenc}
\usepackage[T1]{fontenc}
\usepackage{graphicx}
\usepackage{longtable}
\usepackage{wrapfig}
\usepackage{rotating}
\usepackage[normalem]{ulem}
\usepackage{amsmath}
\usepackage{amssymb}
\usepackage{capt-of}
\usepackage{hyperref}
\author{Cisco Ramon}
\date{\textit{{[}2026-03-14 Sat 00:35]}}
\title{ClusterSetup}
\hypersetup{
 pdfauthor={Cisco Ramon},
 pdftitle={ClusterSetup},
 pdfkeywords={},
 pdfsubject={Building a mini HPC cluster environment using Rocky Linux, VirtualBox, Slurm, MPI, and OpenMP},
 pdfcreator={},
 pdflang={English}}
\begin{document}

\maketitle
\setcounter{tocdepth}{2}
\tableofcontents

\section{Abstract}
\label{sec:org1b8916a}

This documentation presents the complete setup and configuration of a
lightweight High Performance Computing (HPC) cluster environment using
Rocky Linux and VirtualBox. The guide covers practical implementation
of core HPC infrastructure components including shared storage, Slurm
workload management, LDAP-based centralized authentication, MPI runtime
environment, and centralized software stack management using Spack.

The cluster environment supports multinode job execution, centralized
software distribution, distributed MPI applications, and scheduler-aware
runtime execution. The setup follows real-world HPC concepts while
remaining suitable for learning, experimentation, and practical hands-on
training.
\section{Preface}
\label{sec:org30868d6}

This documentation was prepared with the assistance of AI tools for:

\begin{itemize}
\item structuring
\item formatting
\item workflow organization
\item command documentation
\item concept explanations
\end{itemize}

However, the complete HPC cluster setup, configuration, testing, and
validation were performed manually by me in a fully working VirtualBox
based cluster environment.

All commands, scripts, workflows, and configurations documented in this
guide were tested successfully during the setup process, including:

\begin{itemize}
\item cluster networking
\item shared storage
\item NFS configuration
\item Slurm workload manager
\item LDAP authentication
\item centralized user management
\item Spack software stack
\item MPI runtime environment
\item multinode execution
\end{itemize}

This guide is intended for:

\begin{itemize}
\item students
\item beginners
\item Linux learners
\item HPC enthusiasts
\item system administrators
\item researchers
\end{itemize}

who want practical hands-on experience in:

\begin{itemize}
\item HPC infrastructure setup
\item cluster administration
\item workload scheduling
\item centralized authentication
\item parallel programming
\item MPI execution
\item software stack management
\end{itemize}

The goal of this documentation is not only to provide commands, but
also to explain the practical concepts behind real HPC environments in a
simple and beginner-friendly manner.
\section{About the Author}
\label{sec:org830866d}

Hi 👋, I'm Abhishek Raj

I work in the field of High Performance Computing (HPC), Linux systems,
parallel programming, and scientific computing with experience in HPC
application optimization and distributed computing environments,
including the PARAM series of supercomputers.

My primary areas of interest include:

\begin{itemize}
\item HPC clusters and infrastructure
\item Linux systems and administration
\item MPI, OpenMP, CUDA, and OpenACC
\item scientific and research workloads
\item distributed systems
\item software optimization and performance tuning
\item automation and reproducible environments
\end{itemize}

This project was created to gain a deeper understanding of how HPC
clusters are designed and managed internally. The objective was not only
to build a working cluster environment from scratch, but also to better
understand the complete software and infrastructure stack behind
parallel computing systems, which ultimately helps in writing and
optimizing efficient HPC applications.

GitHub Profile:

\textbf{\href{https://github.com/CISSSCO/}{https://github.com/CISSSCO/}}

LinkedIn :

\textbf{\href{https://www.linkedin.com/in/abhi581b}{https://www.linkedin.com/in/abhi581b}}

Project Repository:

\textbf{\href{https://github.com/CISSSCO/hpccfs}{https://github.com/CISSSCO/hpccfs}}

Hosted Documentation:

\textbf{\href{https://cisssco.github.io/hpccfs}{https://cisssco.github.io/hpccfs}}

Source codes, scripts, configuration files, and examples used in this
guide are available in the project repository.

For queries, suggestions, or discussions:

\textbf{\href{https://ciscoramon.pages.dev/}{https://ciscoramon.pages.dev/}}

If you encounter any issue in the documentation, setup process, or
configuration steps, feel free to create an issue in the GitHub
repository.
\section{Introduction}
\label{sec:org6dc219a}

High Performance Computing (HPC) clusters are distributed computing
systems designed to execute computational workloads across multiple
machines simultaneously. Instead of relying on a single system, HPC
clusters combine multiple compute nodes to provide higher computational
capacity, scalability, and parallel execution capability.

Modern HPC environments are widely used in:

\begin{itemize}
\item scientific simulations
\item engineering workloads
\item computational chemistry
\item artificial intelligence
\item machine learning
\item weather modeling
\item large-scale data analysis
\end{itemize}

A typical HPC cluster contains several important components working
together:

\begin{center}
\begin{tabular}{ll}
Component & Purpose\\
\hline
Login Node & User access and job submission\\
Compute Nodes & Execute computational workloads\\
Shared Storage & Shared user data and software\\
Scheduler & Resource allocation and job management\\
MPI Runtime & Distributed parallel execution\\
Authentication System & Centralized user management\\
\end{tabular}
\end{center}

This guide focuses on building these components step-by-step in a small
virtualized environment using:

\begin{itemize}
\item Rocky Linux
\item VirtualBox
\item Slurm
\item OpenMPI
\item LDAP
\item NFS
\item Spack
\end{itemize}

Although the setup is lightweight and designed for learning purposes,
the overall architecture closely follows concepts used in production HPC
centers and supercomputers.

Throughout this documentation, readers will gradually configure:

\begin{itemize}
\item cluster networking
\item shared storage
\item centralized authentication
\item workload scheduling
\item distributed software stacks
\item MPI-based parallel execution
\item scheduler-aware job execution
\end{itemize}

The final environment provides a fully working multinode HPC cluster
capable of running distributed parallel applications across multiple
compute nodes.
\section{Cluster Architecture}
\label{sec:org3975cc6}

The cluster consists of five virtual machines connected using an internal
network.

\begin{verbatim}
                 HOST MACHINE
                       |
         ssh -p XXXX localhost
                       |
               [ NAT Adapter ]
                       |
 ------------------------------------------------
 |               INTERNAL NETWORK               |
 |                  "hpcnet"                    |
 ------------------------------------------------

master      10.10.10.1
login       10.10.10.2
compute01   10.10.10.11
compute02   10.10.10.12
compute03   10.10.10.13
\end{verbatim}
\section{Node Details}
\label{sec:org7cd68ab}
\subsection{Node Role}
\label{sec:org360a002}
\begin{enumerate}
\item Master Node
\label{sec:orga0e0db0}

The master node manages the overall cluster.

Main responsibilities:

\begin{itemize}
\item Slurm controller
\item NFS server
\item LDAP server
\item Monitoring tools
\item Shared software stack
\item Cluster configuration
\end{itemize}
\item Login Node
\label{sec:org156a8df}

The login node is the access point for users.

Typical tasks performed here:

\begin{itemize}
\item SSH login
\item Code compilation
\item Job submission
\item Accessing shared storage
\end{itemize}
\item Compute Nodes
\label{sec:org5b286db}

Compute nodes run the actual workloads.

These nodes are mainly used for:

\begin{itemize}
\item MPI jobs
\item OpenMP jobs
\item Batch workloads
\item Parallel execution
\end{itemize}
\end{enumerate}
\subsection{Networking Setup}
\label{sec:org7c87208}

Two network adapters are used for each VM.
\begin{enumerate}
\item NAT Adapter
\label{sec:orgbbc155e}

Used for:

\begin{itemize}
\item Internet access
\item Package installation
\item SSH access from host system
\end{itemize}
\item Internal Network
\label{sec:orga7b6543}

Used for communication between cluster nodes.

Network name:

\begin{verbatim}
hpcnet
\end{verbatim}

This network is isolated and only used inside the cluster.
\end{enumerate}
\subsection{Static IP Addressing}
\label{sec:orga12700a}

Each node uses a fixed IP address.

\begin{center}
\begin{tabular}{llrr}
VM Name & Hostname & Internal IP & SSH Port\\
\hline
master & master.local & 10.10.10.1 & 2221\\
login & login.local & 10.10.10.2 & 2222\\
compute01 & compute01.local & 10.10.10.11 & 2223\\
compute02 & compute02.local & 10.10.10.12 & 2224\\
compute03 & compute03.local & 10.10.10.13 & 2225\\
\end{tabular}
\end{center}

Example SSH access:

\begin{verbatim}
ssh -p 2221 hpcuser@localhost
\end{verbatim}
\section{Resource Allocation}
\label{sec:org88dc6ff}

The following configuration is enough for a lightweight HPC lab.
You can increase resources later depending on your host system.
\subsection{Master Node}
\label{sec:org0fba0cc}

\begin{center}
\begin{tabular}{ll}
Resource & Value\\
\hline
RAM & 2 GB\\
CPU & 2\\
Disk & 30 GB\\
\end{tabular}
\end{center}
\subsection{Login Node}
\label{sec:org31941a8}

\begin{center}
\begin{tabular}{ll}
Resource & Value\\
\hline
RAM & 2 GB\\
CPU & 2\\
Disk & 25 GB\\
\end{tabular}
\end{center}
\subsection{Compute Nodes}
\label{sec:orgeead451}

\begin{center}
\begin{tabular}{ll}
Resource & Value\\
\hline
RAM & 2 GB\\
CPU & 2\\
Disk & 20 GB\\
\end{tabular}
\end{center}
\section{Master Node Setup}
\label{sec:orgad89bcc}

This section covers the complete setup of the master node. The master node
acts as the central management node of the cluster.

Main responsibilities:

\begin{itemize}
\item Cluster management
\item Shared storage
\item Slurm controller
\item Monitoring services
\item User authentication
\item Software management
\end{itemize}
\subsection{Install VirtualBox on Manjaro}
\label{sec:orga953f9f}

I'm using Manjaro (Arch) as my base system.
\begin{enumerate}
\item Update packages
\label{sec:org103ca1d}

\begin{verbatim}
sudo pacman -Syu
\end{verbatim}
\item Install VirtualBox and host modules
\label{sec:orgc2fa9a2}

\begin{verbatim}
sudo pacman -S virtualbox virtualbox-host-modules-arch
\end{verbatim}
\item Install VirtualBox extension pack
\label{sec:org11e2f0f}

\begin{verbatim}
sudo pacman -S virtualbox-ext-oracle
\end{verbatim}
\item Load VirtualBox kernel module
\label{sec:orgc7b496e}

\begin{verbatim}
sudo modprobe vboxdrv
\end{verbatim}
\item Add current user to VirtualBox group
\label{sec:org3e67ddc}

\begin{verbatim}
sudo usermod -aG vboxusers $USER
\end{verbatim}
\item Reboot the system
\label{sec:orgfe8fcce}

\begin{verbatim}
sudo reboot
\end{verbatim}
\end{enumerate}
\subsection{Download Rocky Linux}
\label{sec:orgc22ca97}

Download Rocky Linux Minimal ISO:

\url{https://rockylinux.org/download}

Choose:

\begin{itemize}
\item x86\_64
\item Minimal ISO
\end{itemize}

The minimal installation keeps the environment lightweight and avoids
unnecessary packages.
\subsection{Create Master VM}
\label{sec:org526e136}

Open VirtualBox and create a new virtual machine.
\begin{enumerate}
\item VM Name
\label{sec:org616298a}

\begin{verbatim}
master
\end{verbatim}
\item Type
\label{sec:org92673d4}

\begin{verbatim}
Linux
Red Hat (64-bit)
\end{verbatim}
\item Hardware
\label{sec:org4f9ae26}

\begin{center}
\begin{tabular}{ll}
Setting & Value\\
\hline
RAM & 2046 MB\\
CPU & 2\\
Disk & 30 GB\\
\end{tabular}
\end{center}
\item Disk Type
\label{sec:org5e2c7ec}

\begin{verbatim}
VDI
Dynamically Allocated
\end{verbatim}

This configuration is enough for a lightweight HPC lab environment.
\end{enumerate}
\subsection{Network Configuration}
\label{sec:org002d291}
\begin{enumerate}
\item Adapter 1
\label{sec:org4bdddb8}

\begin{verbatim}
Attached to: NAT
\end{verbatim}
Used for:

\begin{itemize}
\item Internet access
\item Package installation
\item SSH access from host machine
\end{itemize}
\item NAT Port Forwarding
\label{sec:orgfe9afc7}

Open:

\begin{verbatim}
Settings -> Network -> Adapter 1 -> Advanced -> Port Forwarding
\end{verbatim}

\begin{center}
\begin{tabular}{llrr}
Name & Protocol & Host Port & Guest Port\\
\hline
ssh & TCP & 2221 & 22\\
\end{tabular}
\end{center}
\item Adapter 2
\label{sec:orga4a050a}

\begin{verbatim}
Attached to: Internal Network
\end{verbatim}

Internal Network Name:

\begin{verbatim}
hpcnet
\end{verbatim}

IMPORTANT:
All nodes MUST use the same internal network name.
This network is used for communication between cluster nodes.
\end{enumerate}
\subsection{Rocky Linux Installation}
\label{sec:org5c0d95a}

Choose:

\begin{verbatim}
Minimal Install
\end{verbatim}
\begin{enumerate}
\item Network and Hostname
\label{sec:org0f9c6ac}

Enable BOTH network adapters.

Set hostname:

\begin{verbatim}
master.local
\end{verbatim}
\item User Creation
\label{sec:org9fb4ff5}

Create an administrative user during installation.

\begin{center}
\begin{tabular}{ll}
Field & Value\\
\hline
Username & hpcadmin\\
Administrator & YES\\
\end{tabular}
\end{center}
After installation, log into the VM.
\end{enumerate}
\subsection{Initial System Verification}
\label{sec:org85d0177}
\begin{enumerate}
\item Check Interfaces
\label{sec:orgdbdd2b6}

\begin{verbatim}
ip addr
\end{verbatim}

Expected:

\begin{itemize}
\item NAT interface should receive:
\begin{itemize}
\item 10.0.2.15
\end{itemize}
\item Internal interface should exist but may not yet have static IP
\end{itemize}
\end{enumerate}
\subsection{Configure Static Internal Network}
\label{sec:org2c4db32}
\begin{enumerate}
\item Check Device Names
\label{sec:org22d4d86}

\begin{verbatim}
nmcli device status
\end{verbatim}

Expected:

\begin{center}
\begin{tabular}{ll}
Interface & Purpose\\
\hline
enp0s3 & NAT\\
enp0s8 & Internal Network\\
\end{tabular}
\end{center}
\end{enumerate}
\subsection{Configure Internal Interface}
\label{sec:org64fa059}

Launch NetworkManager TUI:

\begin{verbatim}
sudo nmtui
\end{verbatim}

Navigate:

\begin{verbatim}
Edit a connection
\end{verbatim}

Select:

\begin{verbatim}
enp0s8
\end{verbatim}

Set:

\begin{center}
\begin{tabular}{ll}
Field & Value\\
\hline
IPv4 Method & Manual\\
Address & 10.10.10.1/24\\
Gateway & blank\\
DNS & blank\\
\end{tabular}
\end{center}

DO NOT modify NAT interface.
\subsection{Restart Networking}
\label{sec:org01285b9}

\begin{verbatim}
sudo systemctl restart NetworkManager
\end{verbatim}
\subsection{Verify Networking}
\label{sec:org6f55435}

\begin{verbatim}
ip addr
\end{verbatim}

Expected:

\begin{center}
\begin{tabular}{lr}
Interface & IP\\
\hline
NAT & 10.0.2.x\\
Internal & 10.10.10.1\\
\end{tabular}
\end{center}
\subsection{SSH Access from Host}
\label{sec:org9569b21}

From host machine:

\begin{verbatim}
ssh -p 2221 hpcadmin@localhost
\end{verbatim}
If the connection works, port forwarding is configured correctly.
\subsection{System Update and Packages}
\label{sec:orge075d3a}
\begin{enumerate}
\item Update System
\label{sec:orgca9c7fb}

\begin{verbatim}
sudo dnf update -y
\end{verbatim}
\item Install Base Packages
\label{sec:org1973a59}

\begin{verbatim}
sudo dnf install -y \
vim wget curl git \
net-tools bind-utils \
htop tmux tree \
gcc gcc-c++ gcc-gfortran \
make cmake \
python3 \
openssh-server
\end{verbatim}

These packages provide:

\begin{itemize}
\item Development tools
\item Network utilities
\item Monitoring utilities
\item Basic administration tools
\end{itemize}
\end{enumerate}
\subsection{Timezone Configuration}
\label{sec:org178e2cc}

Check timezone:

\begin{verbatim}
timedatectl
\end{verbatim}

Set timezone if needed:

\begin{verbatim}
sudo timedatectl set-timezone Asia/Kolkata
\end{verbatim}
\subsection{Configure /etc/hosts}
\label{sec:org32af993}

Edit:

\begin{verbatim}
sudo vim /etc/hosts
\end{verbatim}

Add:

\begin{verbatim}
10.10.10.1   master master.local
10.10.10.2   login login.local
10.10.10.11  compute01 compute01.local
10.10.10.12  compute02 compute02.local
10.10.10.13  compute03 compute03.local
\end{verbatim}
This allows hostname-based communication between nodes.
\subsection{Disable Firewall (Temporary)}
\label{sec:orgf01d714}

For initial setup and testing:
\begin{verbatim}
sudo systemctl disable --now firewalld
\end{verbatim}
NOTE:

Firewall rules can be configured properly later after cluster setup is complete.
\subsection{SELinux Configuration}
\label{sec:orgaf1bbf6}

Check status:

\begin{verbatim}
getenforce
\end{verbatim}

Temporarily disable enforcing:

\begin{verbatim}
sudo setenforce 0
\end{verbatim}

Persist configuration:

\begin{verbatim}
sudo vim /etc/selinux/config
\end{verbatim}

Modify:

\begin{verbatim}
SELINUX=permissive
\end{verbatim}
This helps avoid permission-related issues during initial cluster setup.
\subsection{Verify Hostname}
\label{sec:org3379bde}

\begin{verbatim}
hostnamectl
\end{verbatim}

Expected:

\begin{verbatim}
Static hostname: master.local
\end{verbatim}

If hostname is not updated(showing previous), reboot the machine.

At this point, the master node setup is complete and ready for cloning or
further cluster configuration.
\section{Cloning Nodes}
\label{sec:org733443a}

Purpose:

\begin{itemize}
\item avoid reinstalling all nodes manually
\item simulate real HPC image provisioning workflow
\end{itemize}
\subsection{Create Snapshot}
\label{sec:org5973cea}

Power off VM:

\begin{verbatim}
sudo poweroff
\end{verbatim}

Create snapshot in VirtualBox:

\begin{verbatim}
clean-base-master
\end{verbatim}
\subsection{Clone Nodes}
\label{sec:org76ac096}

Create FULL CLONES.

IMPORTANT:
Enable:

\begin{verbatim}
Reinitialize the MAC address of all network cards
\end{verbatim}
\subsection{Clone Names}
\label{sec:org95359c9}

\begin{center}
\begin{tabular}{l}
VM Name\\
\hline
login\\
compute01\\
compute02\\
compute03\\
\end{tabular}
\end{center}
\subsection{Configure Port Forwarding}
\label{sec:org8faf02a}
\begin{enumerate}
\item Login Node
\label{sec:org9f40653}

\begin{center}
\begin{tabular}{rr}
Host Port & Guest Port\\
\hline
2222 & 22\\
\end{tabular}
\end{center}
\item Compute01
\label{sec:org2f41b03}

\begin{center}
\begin{tabular}{rr}
Host Port & Guest Port\\
\hline
2223 & 22\\
\end{tabular}
\end{center}
\item Compute02
\label{sec:org7731d76}

\begin{center}
\begin{tabular}{rr}
Host Port & Guest Port\\
\hline
2224 & 22\\
\end{tabular}
\end{center}
\item Compute03
\label{sec:org6a58f77}

\begin{center}
\begin{tabular}{rr}
Host Port & Guest Port\\
\hline
2225 & 22\\
\end{tabular}
\end{center}
\end{enumerate}
\subsection{Machine ID Regeneration}
\label{sec:org0f71801}

IMPORTANT:
Each cloned Linux machine must have unique machine-id.
\subsection{Remove Existing Machine ID}
\label{sec:org24d3f81}

\begin{verbatim}
sudo rm -f /etc/machine-id
\end{verbatim}
\subsection{Generate New Machine ID}
\label{sec:orgc4669f6}

\begin{verbatim}
sudo systemd-machine-id-setup
\end{verbatim}
\subsection{Verify Machine ID}
\label{sec:org2f4837c}

\begin{verbatim}
cat /etc/machine-id
\end{verbatim}

Ensure every node has unique machine-id.
\subsection{Configure Hostname}
\label{sec:orgabdd5e0}
\begin{enumerate}
\item Login Node
\label{sec:org8bbf406}

Edit hostname:

\begin{verbatim}
sudo vim /etc/hostname
\end{verbatim}

Set:

\begin{verbatim}
login.local
\end{verbatim}

Apply:

\begin{verbatim}
sudo hostnamectl set-hostname login.local
\end{verbatim}
\item Compute01
\label{sec:orgd181b27}

Set hostname:

\begin{verbatim}
compute01.local
\end{verbatim}

Apply:

\begin{verbatim}
sudo hostnamectl set-hostname compute01.local
\end{verbatim}
\item Compute02
\label{sec:orgabf4e39}

\begin{verbatim}
compute02.local
\end{verbatim}

\begin{verbatim}
sudo hostnamectl set-hostname compute02.local
\end{verbatim}
\item Compute03
\label{sec:org2d145b2}

\begin{verbatim}
compute03.local
\end{verbatim}

\begin{verbatim}
sudo hostnamectl set-hostname compute03.local
\end{verbatim}
\end{enumerate}
\subsection{Configure Static Internal IPs}
\label{sec:org3c2895d}

Launch:

\begin{verbatim}
sudo nmtui
\end{verbatim}

Edit internal adapter:

\begin{verbatim}
enp0s8
\end{verbatim}
\begin{enumerate}
\item Login Node
\label{sec:org300ecde}

\begin{verbatim}
10.10.10.2/24
\end{verbatim}
\item Compute01
\label{sec:orga542e6e}

\begin{verbatim}
10.10.10.11/24
\end{verbatim}
\item Compute02
\label{sec:org8086f38}

\begin{verbatim}
10.10.10.12/24
\end{verbatim}
\item Compute03
\label{sec:orgc87057f}

\begin{verbatim}
10.10.10.13/24
\end{verbatim}
\end{enumerate}
\subsection{Restart Networking}
\label{sec:org2467d23}

\begin{verbatim}
sudo systemctl restart NetworkManager
\end{verbatim}
\subsection{Verify Networking}
\label{sec:orgb6b358d}

\begin{verbatim}
ip addr
\end{verbatim}
\subsection{Verify Routing}
\label{sec:orgc5c36ad}

\begin{verbatim}
ip route
\end{verbatim}

Expected:

\begin{verbatim}
default via 10.0.2.2 dev enp0s3
10.10.10.0/24 dev enp0s8
\end{verbatim}
\subsection{Verify Connectivity}
\label{sec:org99bb80a}
\begin{enumerate}
\item Ping Between Nodes
\label{sec:org3845f66}

\begin{verbatim}
ping master
\end{verbatim}

\begin{verbatim}
ping login
\end{verbatim}

\begin{verbatim}
ping compute01
\end{verbatim}
\end{enumerate}
\subsection{SSH Verification}
\label{sec:org16363e9}

From host:
\begin{enumerate}
\item Master
\label{sec:org5e1b9ac}

\begin{verbatim}
ssh -p 2221 hpcadmin@localhost
\end{verbatim}
\item Login
\label{sec:org8ea4296}

\begin{verbatim}
ssh -p 2222 hpcadmin@localhost
\end{verbatim}
\item Compute01
\label{sec:orgb5bfbfd}

\begin{verbatim}
ssh -p 2223 hpcadmin@localhost
\end{verbatim}
\item Compute02
\label{sec:org92b01ee}

\begin{verbatim}
ssh -p 2224 hpcadmin@localhost
\end{verbatim}
\item Compute03
\label{sec:org2535fc9}

\begin{verbatim}
ssh -p 2225 hpcadmin@localhost
\end{verbatim}
\end{enumerate}
\section{Shared Infrastructure Services}
\label{sec:orgc36def5}

Goal:

Transform the cluster into a realistic HPC-style shared environment.

This phase includes:

\begin{itemize}
\item passwordless SSH
\item shared home directories
\item NFS server/client configuration
\item persistent mounts
\item multi-node shared filesystem verification
\end{itemize}

By the end of this phase:

\begin{itemize}
\item all nodes share same /home
\item users can access same files from all nodes
\item cluster behaves like a real HPC environment
\end{itemize}
\subsection{Cluster Role Architecture}
\label{sec:orgb54ea02}
\begin{enumerate}
\item Master Node
\label{sec:orgb616ec6}

Purpose:

\begin{itemize}
\item cluster management
\item NFS server
\item future Slurm controller
\item future Munge authority
\item future LDAP server
\end{itemize}

Services:

\begin{center}
\begin{tabular}{ll}
Service & Purpose\\
\hline
nfs-server & shared storage\\
future slurmctld & scheduler controller\\
future munge & authentication\\
\end{tabular}
\end{center}
\item Login Node
\label{sec:org2175bb6}

Purpose:

\begin{itemize}
\item user login
\item compilation
\item job submission
\end{itemize}

Users should:
\begin{itemize}
\item SSH into login node
\item compile programs
\item submit jobs using Slurm
\end{itemize}
\item Compute Nodes
\label{sec:orgc3a174b}

Purpose:

\begin{itemize}
\item execute jobs only
\end{itemize}

Future services:

\begin{center}
\begin{tabular}{ll}
Service & Purpose\\
\hline
slurmd & compute daemon\\
munge & authentication\\
\end{tabular}
\end{center}
\end{enumerate}
\subsection{Verify User Consistency}
\label{sec:org664b62a}

IMPORTANT:

NFS shared storage requires same UID/GID across all nodes.

Verify user identity on all nodes:

\begin{verbatim}
id hpcadmin
\end{verbatim}

Expected:

\begin{verbatim}
uid=1000(hpcadmin) gid=1000(hpcadmin)
\end{verbatim}

All nodes must match.
\subsection{Passwordless SSH Configuration}
\label{sec:org71cbb1c}

Purpose:

\begin{itemize}
\item MPI communication
\item Slurm execution
\item cluster automation
\item remote management
\end{itemize}

SSH trust is generated ONLY from login node.

This mimics real HPC architecture:

\begin{verbatim}
User -> Login Node -> Compute Nodes
\end{verbatim}
\subsection{Generate SSH Key on Login Node}
\label{sec:orge43bf8f}

SSH into login node from host:

\begin{verbatim}
ssh -p 2222 hpcadmin@localhost
\end{verbatim}

Generate SSH key:

\begin{verbatim}
ssh-keygen
\end{verbatim}

Press Enter for all prompts.

Generated files:

\begin{verbatim}
~/.ssh/id_rsa
~/.ssh/id_rsa.pub
\end{verbatim}
\subsection{Copy SSH Keys to Cluster Nodes}
\label{sec:orgc25a067}
\begin{enumerate}
\item Master
\label{sec:orge5b5e4d}

\begin{verbatim}
ssh-copy-id hpcadmin@master
\end{verbatim}
\item Compute01
\label{sec:orgc59a6e6}

\begin{verbatim}
ssh-copy-id hpcadmin@compute01
\end{verbatim}
\item Compute02
\label{sec:org7afd088}

\begin{verbatim}
ssh-copy-id hpcadmin@compute02
\end{verbatim}
\item Compute03
\label{sec:org7f249e6}

\begin{verbatim}
ssh-copy-id hpcadmin@compute03
\end{verbatim}
\end{enumerate}
\subsection{Verify Passwordless SSH}
\label{sec:org02af0f7}
\begin{enumerate}
\item Test Master
\label{sec:orge16c6b4}

\begin{verbatim}
ssh master hostname
\end{verbatim}

Expected:

\begin{verbatim}
master.local
\end{verbatim}
\item Test Compute01
\label{sec:org6598b89}

\begin{verbatim}
ssh compute01 hostname
\end{verbatim}

Expected:

\begin{verbatim}
compute01.local
\end{verbatim}
\item Test Compute02
\label{sec:org089e802}

\begin{verbatim}
ssh compute02 hostname
\end{verbatim}
\item Test Compute03
\label{sec:orgf23e008}

\begin{verbatim}
ssh compute03 hostname
\end{verbatim}

No password prompts should appear.
\end{enumerate}
\subsection{HPC Concept — Why Passwordless SSH Matters}
\label{sec:org24c6090}

Passwordless SSH becomes essential for:

\begin{center}
\begin{tabular}{ll}
Feature & Purpose\\
\hline
MPI & remote process launch\\
Slurm & distributed execution\\
automation & orchestration\\
monitoring & remote administration\\
\end{tabular}
\end{center}
\subsection{NFS Shared Storage Setup}
\label{sec:org8864621}

Purpose:

Provide shared home directories across all cluster nodes.

This is a core HPC architecture requirement.

Without shared storage:

\begin{itemize}
\item source code differs per node
\item binaries differ
\item scripts unavailable across nodes
\item MPI execution becomes difficult
\end{itemize}

Real HPC systems commonly use:

\begin{itemize}
\item NFS
\item Lustre
\item BeeGFS
\item GPFS
\end{itemize}

This lab begins with NFS.
\subsection{Install NFS Packages}
\label{sec:orge985083}

Install on ALL nodes:

\begin{verbatim}
sudo dnf install -y nfs-utils
\end{verbatim}
\subsection{Configure NFS Server on Master}
\label{sec:orge404015}

SSH into master node.
\subsection{Enable NFS Services}
\label{sec:org8398431}

\begin{verbatim}
sudo systemctl enable --now nfs-server
\end{verbatim}

Verify:

\begin{verbatim}
systemctl status nfs-server
\end{verbatim}

Expected:

\begin{verbatim}
active (running)
\end{verbatim}
\subsection{Configure NFS Exports}
\label{sec:org33dc423}

Edit exports file:

\begin{verbatim}
sudo vim /etc/exports
\end{verbatim}

Add:

\begin{verbatim}
/home 10.10.10.0/24(rw,sync,no_root_squash)
\end{verbatim}
\subsection{Export Option Explanation}
\label{sec:org83df700}

\begin{center}
\begin{tabular}{ll}
Option & Meaning\\
\hline
rw & read/write access\\
sync & synchronous writes\\
no\_root\_squash & preserve root privileges\\
\end{tabular}
\end{center}
\subsection{Apply Exports}
\label{sec:org005e15a}

\begin{verbatim}
sudo exportfs -rav
\end{verbatim}
\subsection{Verify Exports}
\label{sec:org06ed5b5}

\begin{verbatim}
sudo exportfs -v
\end{verbatim}

Expected export:

\begin{verbatim}
/home 10.10.10.0/24(...)
\end{verbatim}
\subsection{Configure NFS Clients}
\label{sec:orgc2539e6}

Perform on:

\begin{itemize}
\item login
\item compute01
\item compute02
\item compute03
\end{itemize}
\subsection{Backup Existing Home Directory}
\label{sec:org44a30a3}

\begin{verbatim}
sudo mv /home /home.local
\end{verbatim}

Purpose:

Preserve original local home directory.
\subsection{Create New Mount Point}
\label{sec:orgab6dcea}

\begin{verbatim}
sudo mkdir /home
\end{verbatim}
\subsection{Mount Shared NFS Home}
\label{sec:org923778d}

\begin{verbatim}
sudo mount master:/home /home
\end{verbatim}
\subsection{Verify Mount}
\label{sec:org086a3ec}

\begin{verbatim}
df -h
\end{verbatim}

Expected:

\begin{verbatim}
master:/home
\end{verbatim}

mounted on:

\begin{verbatim}
/home
\end{verbatim}
\subsection{Verify Shared Filesystem}
\label{sec:orgdcb3fa9}

From login node:

\begin{verbatim}
touch ~/shared_test
\end{verbatim}

From compute01:

\begin{verbatim}
ls ~
\end{verbatim}

Expected:

\begin{verbatim}
shared_test
\end{verbatim}

This confirms successful shared filesystem operation.
\subsection{Configure Persistent NFS Mounts}
\label{sec:orgf32234c}

Without this configuration:

\begin{itemize}
\item NFS mount disappears after reboot
\end{itemize}
\subsection{Configure /etc/fstab}
\label{sec:orgfc03d41}

On ALL client nodes:

\begin{verbatim}
sudo vim /etc/fstab
\end{verbatim}

Add:

\begin{verbatim}
master:/home   /home   nfs   defaults   0 0
\end{verbatim}
\subsection{Verify fstab Before Reboot}
\label{sec:org5ef16db}

IMPORTANT:

Always validate fstab before rebooting.

Run:

\begin{verbatim}
sudo mount -a
\end{verbatim}

If no errors appear:
\begin{itemize}
\item fstab configuration is correct
\end{itemize}
\subsection{Reboot Validation Procedure}
\label{sec:org38acde6}

Purpose:

Validate persistent cluster infrastructure after reboot.

This is important real-world HPC administration practice.
\subsection{Reboot Order}
\label{sec:orgd14e8be}
\begin{enumerate}
\item Reboot Master First
\label{sec:org5bc1b4e}

\begin{verbatim}
sudo reboot
\end{verbatim}

Wait until:
\begin{itemize}
\item SSH becomes available
\item NFS server active
\end{itemize}

Verify:

\begin{verbatim}
systemctl status nfs-server
\end{verbatim}
\item Reboot Login Node
\label{sec:org87e0749}

\begin{verbatim}
sudo reboot
\end{verbatim}
\item Reboot Compute Nodes
\label{sec:org0aea1f8}

\begin{verbatim}
sudo reboot
\end{verbatim}

Perform for:
\begin{itemize}
\item compute01
\item compute02
\item compute03
\end{itemize}
\end{enumerate}
\subsection{Verify Persistent Mounts After Reboot}
\label{sec:orgaaa11ff}

On login node:

\begin{verbatim}
df -h
\end{verbatim}

Expected:

\begin{verbatim}
master:/home
\end{verbatim}

mounted automatically.
\subsection{Verify Active NFS Mount}
\label{sec:org235be98}

\begin{verbatim}
mount | grep home
\end{verbatim}

Expected:

\begin{verbatim}
master:/home on /home type nfs
\end{verbatim}
\subsection{Verify Shared Files After Reboot}
\label{sec:org37b7523}

\begin{verbatim}
ls ~
\end{verbatim}

Expected:
\begin{itemize}
\item previously created shared files still visible
\end{itemize}
\section{Munge Authentication + Slurm Scheduler}
\label{sec:org3171f44}

Goal:

Transform the cluster into a scheduler-managed HPC environment.

This phase includes:

\begin{itemize}
\item Munge authentication
\item Slurm controller setup
\item Compute node daemon setup
\item Scheduler configuration
\item Partition creation
\item Multi-node execution
\item Resource allocation testing
\end{itemize}

By the end of this phase:

\begin{itemize}
\item cluster nodes authenticate securely
\item Slurm manages compute resources
\item jobs execute across nodes
\item scheduler topology becomes operational
\end{itemize}
\subsection{HPC Scheduler Architecture}
\label{sec:org49d3137}
\begin{enumerate}
\item Master Node
\label{sec:orgcff816a}

Purpose:

\begin{itemize}
\item Slurm controller
\item cluster scheduler
\item resource manager
\item Munge authority
\end{itemize}

Services:

\begin{center}
\begin{tabular}{ll}
Service & Purpose\\
\hline
munged & cluster authentication\\
slurmctld & scheduler controller\\
\end{tabular}
\end{center}
\item Login Node
\label{sec:orgd641f5c}

Purpose:

\begin{itemize}
\item user login
\item compilation
\item job submission
\end{itemize}

Users should:

\begin{itemize}
\item login here
\item compile here
\item submit jobs here
\end{itemize}
\item Compute Nodes
\label{sec:org1f30fda}

Purpose:

\begin{itemize}
\item execute jobs
\end{itemize}

Services:

\begin{center}
\begin{tabular}{ll}
Service & Purpose\\
\hline
munged & authentication\\
slurmd & compute execution daemon\\
\end{tabular}
\end{center}
\end{enumerate}
\subsection{IMPORTANT Rocky Linux 10 Issue}
\label{sec:org2be00f1}

Originally the cluster was created using:

\begin{verbatim}
Rocky Linux 10
\end{verbatim}

Problem encountered:

\begin{verbatim}
No match for argument: slurm
\end{verbatim}

Reason:

\begin{itemize}
\item Rocky Linux 10 HPC ecosystem still immature
\item Slurm packages unavailable in standard repos
\item OpenHPC EL10 support incomplete
\end{itemize}

Resolution:

\begin{itemize}
\item cluster rebuilt using Rocky Linux 9 Minimal
\end{itemize}

Important HPC lesson:

Production HPC clusters usually avoid bleeding-edge enterprise releases because:

\begin{itemize}
\item scheduler ecosystems lag behind
\item HPC repos may not exist
\item compatibility becomes unstable
\end{itemize}
\subsection{Install EPEL Repository}
\label{sec:orgbba9833}

Install on ALL nodes:

\begin{verbatim}
sudo dnf install -y epel-release
\end{verbatim}
\subsection{Enable CRB Repository}
\label{sec:org1f3685e}

Required on Rocky Linux 9 and 10.

Install on ALL nodes:

\begin{verbatim}
sudo crb enable
\end{verbatim}
\subsection{Install Munge + Slurm Packages}
\label{sec:orgd046c3d}

Install on ALL nodes:

\begin{verbatim}
sudo dnf install -y \
munge munge-libs \
slurm slurm-slurmd slurm-slurmctld
\end{verbatim}
\subsection{Verify Installed Packages}
\label{sec:org53b1698}
\begin{enumerate}
\item Verify Munge
\label{sec:orgec1ad7c}

\begin{verbatim}
rpm -qa | grep munge
\end{verbatim}

Expected:

\begin{verbatim}
munge
munge-libs
\end{verbatim}
\item Verify Slurm
\label{sec:orgd41c25e}

\begin{verbatim}
rpm -qa | grep slurm
\end{verbatim}

Expected:

\begin{verbatim}
slurm
slurm-slurmd
slurm-slurmctld
\end{verbatim}
\end{enumerate}
\subsection{Verify System Users}
\label{sec:orgbd4263d}
\begin{enumerate}
\item Verify Munge User
\label{sec:org98092b4}

\begin{verbatim}
id munge
\end{verbatim}
\item Verify Slurm User
\label{sec:org1fcc054}

\begin{verbatim}
id slurm
\end{verbatim}
\end{enumerate}
\subsection{IMPORTANT Issue Encountered — Missing slurm User}
\label{sec:org99b1c8b}

Problem:

\begin{verbatim}
id: 'slurm': no such user
\end{verbatim}

Reason:

Some Slurm packages on Rocky Linux do not automatically create the slurm user.

Resolution:

Create manually on ALL nodes.
\subsection{Create Slurm Group}
\label{sec:orgdbd4a2c}

\begin{verbatim}
sudo groupadd slurm
\end{verbatim}
\subsection{Create Slurm User}
\label{sec:orgd096d31}

\begin{verbatim}
sudo useradd -r -g slurm -s /sbin/nologin slurm
\end{verbatim}

Explanation:

\begin{center}
\begin{tabular}{ll}
Option & Meaning\\
\hline
-r & system account\\
-g slurm & primary group\\
-s /sbin/nologin & disable shell login\\
\end{tabular}
\end{center}
\subsection{Verify Slurm User}
\label{sec:org0c67c03}

\begin{verbatim}
id slurm
\end{verbatim}

Expected:

\begin{verbatim}
uid=...
gid=...
\end{verbatim}
\subsection{IMPORTANT HPC Security Concept}
\label{sec:orgab98ae2}

Cluster daemons should not run permanently as root.

Therefore:

\begin{itemize}
\item Munge runs as munge user
\item Slurm runs as slurm user
\end{itemize}

This is standard enterprise daemon security practice.
\subsection{Configure Munge Authentication}
\label{sec:org65d51cb}

Purpose:

Create secure trust relationship between all cluster nodes.

IMPORTANT:

All nodes MUST use the same:

\begin{verbatim}
/etc/munge/munge.key
\end{verbatim}

This key acts as the cluster-wide shared secret.
\subsection{Generate Munge Key on Master}
\label{sec:orga9c8976}

ONLY on master node:

\begin{verbatim}
sudo /usr/sbin/create-munge-key
\end{verbatim}
\subsection{Verify Munge Key}
\label{sec:org0a6ad1e}

\begin{verbatim}
ls -l /etc/munge/munge.key
\end{verbatim}

Expected:

\begin{verbatim}
-rw-------
\end{verbatim}

owned by:

\begin{verbatim}
munge munge
\end{verbatim}
\subsection{IMPORTANT Mistake Encountered — Munge Key Generated on All Nodes}
\label{sec:orgbde064c}

Problem:

Munge key accidentally generated independently on all nodes.

Result:

Nodes could not authenticate each other.

Important HPC lesson:

Munge works using a:

\begin{verbatim}
shared trust model
\end{verbatim}

All nodes MUST share identical munge.key.
\subsection{Correct Munge Key Distribution}
\label{sec:org2b3f27f}
\begin{enumerate}
\item Copy Key Temporarily
\label{sec:orge0fce32}

On master:

\begin{verbatim}
sudo cp /etc/munge/munge.key /tmp/
\end{verbatim}
\item Temporary Permission Adjustment
\label{sec:org2b5726f}

\begin{verbatim}
sudo chmod 644 /tmp/munge.key
\end{verbatim}

Purpose:

Allow normal user SCP transfer.
\end{enumerate}
\subsection{IMPORTANT SSH Root Login Issue}
\label{sec:org88a2235}

Problem encountered:

\begin{verbatim}
Permission denied (publickey,gssapi-keyex,gssapi-with-mic,password)
\end{verbatim}

Reason:

Using:

\begin{verbatim}
sudo scp ...
\end{verbatim}

attempted root SSH login.

Rocky Linux disables direct root SSH login by default.

Resolution:

Use normal user SCP.
\subsection{Copy Munge Key to Login Node}
\label{sec:org25899de}

\begin{verbatim}
scp /tmp/munge.key login:/tmp/
\end{verbatim}
\subsection{Copy Munge Key to Compute01}
\label{sec:org083d945}

\begin{verbatim}
scp /tmp/munge.key compute01:/tmp/
\end{verbatim}
\subsection{Copy Munge Key to Compute02}
\label{sec:org93ad156}

\begin{verbatim}
scp /tmp/munge.key compute02:/tmp/
\end{verbatim}
\subsection{Copy Munge Key to Compute03}
\label{sec:orgaa58fad}

\begin{verbatim}
scp /tmp/munge.key compute03:/tmp/
\end{verbatim}
\subsection{Remove Temporary Key Copy}
\label{sec:org3c1b0b9}

On master:

\begin{verbatim}
sudo rm -f /tmp/munge.key
\end{verbatim}
\subsection{Install Munge Key on Client Nodes}
\label{sec:org1bb49a6}

Perform on:

\begin{itemize}
\item login
\item compute01
\item compute02
\item compute03
\end{itemize}
\subsection{Stop Munge Service}
\label{sec:orgd46f0c4}

\begin{verbatim}
sudo systemctl stop munge
\end{verbatim}
\subsection{Install Key}
\label{sec:orgc4868a0}

\begin{verbatim}
sudo mv /tmp/munge.key /etc/munge/
\end{verbatim}
\subsection{Fix Ownership}
\label{sec:orgafcdc5e}

\begin{verbatim}
sudo chown munge:munge /etc/munge/munge.key
\end{verbatim}
\subsection{Fix Permissions}
\label{sec:org9f17df9}

\begin{verbatim}
sudo chmod 400 /etc/munge/munge.key
\end{verbatim}
\subsection{Start Munge}
\label{sec:orgf31a608}

On ALL nodes:

\begin{verbatim}
sudo systemctl enable --now munge
\end{verbatim}
\subsection{IMPORTANT Issue Encountered — Munge Not Running on Master}
\label{sec:org6a7d65f}

Problem:

\begin{verbatim}
munge: Error: Failed to access "/var/run/munge/munge.socket.2"
\end{verbatim}

Reason:

Munge service never started on master.

Resolution:

\begin{verbatim}
sudo systemctl enable --now munge
\end{verbatim}

Important lesson:

Difference between:

\begin{center}
\begin{tabular}{ll}
State & Meaning\\
\hline
inactive & service not started\\
failed & service crashed\\
\end{tabular}
\end{center}
\subsection{Verify Munge Service}
\label{sec:org8dcdd6a}

\begin{verbatim}
systemctl status munge
\end{verbatim}

Expected:

\begin{verbatim}
active (running)
\end{verbatim}
\subsection{Verify Munge Authentication}
\label{sec:org20fec7f}

From master:

\begin{verbatim}
munge -n | ssh compute01 unmunge
\end{verbatim}
\subsection{Test Compute02}
\label{sec:orgdf93d86}

\begin{verbatim}
munge -n | ssh compute02 unmunge
\end{verbatim}
\subsection{Test Compute03}
\label{sec:org6d7f99c}

\begin{verbatim}
munge -n | ssh compute03 unmunge
\end{verbatim}
\subsection{Test Login Node}
\label{sec:org27aa016}

\begin{verbatim}
munge -n | ssh login unmunge
\end{verbatim}

Expected:

\begin{verbatim}
STATUS: Success
\end{verbatim}
\subsection{HPC Concepts Learned}
\label{sec:orgd56095f}

Munge introduced:

\begin{itemize}
\item cluster-wide authentication
\item distributed trust model
\item shared secret authentication
\item daemon security
\item inter-node trust validation
\end{itemize}
\subsection{Configure Slurm Scheduler}
\label{sec:orga3ca4a1}

Purpose:

Enable centralized resource scheduling.

IMPORTANT:

\begin{verbatim}
/etc/slurm/slurm.conf
\end{verbatim}

must be IDENTICAL on ALL nodes.
\subsection{IMPORTANT Issue Encountered — Default localhost Configuration}
\label{sec:orgdb24485}

Problem:

\begin{verbatim}
sinfo
-> localhost
\end{verbatim}

Reason:

Default example slurm.conf remained active.

Resolution:

Replace configuration with custom cluster topology.

Important HPC lesson:

Incorrect topology configuration is one of the most common Slurm admin problems.
\subsection{Create Slurm Configuration}
\label{sec:org53be14b}

On master:

\begin{verbatim}
sudo vim /etc/slurm/slurm.conf
\end{verbatim}

Replace ENTIRE file with:

\begin{verbatim}
ClusterName=hpccluster
SlurmctldHost=master

MpiDefault=none
ProctrackType=proctrack/cgroup
ReturnToService=2

SlurmctldPidFile=/run/slurmctld.pid
SlurmdPidFile=/run/slurmd.pid

SlurmdSpoolDir=/var/spool/slurmd
StateSaveLocation=/var/spool/slurmctld

SlurmUser=slurm

SwitchType=switch/none
TaskPlugin=task/affinity

SchedulerType=sched/backfill

SelectType=select/cons_tres
SelectTypeParameters=CR_Core

AuthType=auth/munge

NodeName=compute01 CPUs=2 State=UNKNOWN
NodeName=compute02 CPUs=2 State=UNKNOWN
NodeName=compute03 CPUs=2 State=UNKNOWN

PartitionName=debug Nodes=ALL Default=YES MaxTime=INFINITE State=UP
\end{verbatim}
\subsection{Configuration Explanation}
\label{sec:org3ec3395}

\begin{center}
\begin{tabular}{ll}
Parameter & Purpose\\
\hline
ClusterName & cluster identifier\\
SlurmctldHost & controller node\\
NodeName & compute nodes\\
CPUs=2 & available CPUs\\
PartitionName & scheduler queue\\
\end{tabular}
\end{center}
\subsection{Create Slurm Runtime Directories}
\label{sec:orgd4a3e8e}
\begin{enumerate}
\item On Master
\label{sec:org90350b0}

\begin{verbatim}
sudo mkdir -p /var/spool/slurmctld
\end{verbatim}

\begin{verbatim}
sudo mkdir -p /var/spool/slurmd
\end{verbatim}
\end{enumerate}
\subsection{Set Ownership}
\label{sec:org1d34b67}

\begin{verbatim}
sudo chown slurm:slurm /var/spool/slurmctld
\end{verbatim}

\begin{verbatim}
sudo chown slurm:slurm /var/spool/slurmd
\end{verbatim}
\subsection{Create slurmd Directories on Compute Nodes}
\label{sec:org9875168}
\begin{enumerate}
\item Compute01
\label{sec:org51cd6e1}

\begin{verbatim}
ssh compute01 sudo -S mkdir -p /var/spool/slurmd
\end{verbatim}
\item Compute02
\label{sec:org0763254}

\begin{verbatim}
ssh compute02 sudo -S mkdir -p /var/spool/slurmd
\end{verbatim}
\item Compute03
\label{sec:org4d7c703}

\begin{verbatim}
ssh compute03 sudo -S mkdir -p /var/spool/slurmd
\end{verbatim}
\end{enumerate}
\subsection{IMPORTANT sudo Over SSH Issue}
\label{sec:org30393fe}

Problem:

\begin{verbatim}
sudo: a terminal is required to read the password
\end{verbatim}

Resolution:

Use:

\begin{verbatim}
sudo -S
\end{verbatim}

Important Linux administration concept:

sudo over non-interactive SSH requires stdin password handling.
\subsection{Fix Ownership on Compute Nodes}
\label{sec:org5d2a46f}
\begin{enumerate}
\item Compute01
\label{sec:org928b7c4}

\begin{verbatim}
ssh compute01 sudo -S chown slurm:slurm /var/spool/slurmd
\end{verbatim}
\item Compute02
\label{sec:org1fbb616}

\begin{verbatim}
ssh compute02 sudo -S chown slurm:slurm /var/spool/slurmd
\end{verbatim}
\item Compute03
\label{sec:org837ec84}

\begin{verbatim}
ssh compute03 sudo -S chown slurm:slurm /var/spool/slurmd
\end{verbatim}
\end{enumerate}
\subsection{Verify Runtime Directories}
\label{sec:org4941924}

\begin{verbatim}
ls -ld /var/spool/slurmctld
\end{verbatim}

Expected:

\begin{verbatim}
slurm slurm
\end{verbatim}
\subsection{Distribute slurm.conf}
\label{sec:org2fe31d9}
\begin{enumerate}
\item Login
\label{sec:org20eed47}

\begin{verbatim}
scp /etc/slurm/slurm.conf login:/tmp/
\end{verbatim}
\item Compute01
\label{sec:orgbe17bb2}

\begin{verbatim}
scp /etc/slurm/slurm.conf compute01:/tmp/
\end{verbatim}
\item Compute02
\label{sec:org57b3ded}

\begin{verbatim}
scp /etc/slurm/slurm.conf compute02:/tmp/
\end{verbatim}
\item Compute03
\label{sec:org6aad414}

\begin{verbatim}
scp /etc/slurm/slurm.conf compute03:/tmp/
\end{verbatim}
\end{enumerate}
\subsection{Install Configuration on Nodes}
\label{sec:org75e3daa}

On ALL nodes:

\begin{verbatim}
sudo mv /tmp/slurm.conf /etc/slurm/slurm.conf
\end{verbatim}
\subsection{Start Slurm Controller}
\label{sec:org45ba0b1}

ONLY on master:

\begin{verbatim}
sudo systemctl enable --now slurmctld
\end{verbatim}
\subsection{Verify Controller}
\label{sec:org118eea3}

\begin{verbatim}
systemctl status slurmctld
\end{verbatim}

Expected:

\begin{verbatim}
active (running)
\end{verbatim}
\subsection{Start slurmd on Compute Nodes}
\label{sec:org46e32e9}

On:
\begin{itemize}
\item compute01
\item compute02
\item compute03
\end{itemize}

\begin{verbatim}
sudo systemctl enable --now slurmd
\end{verbatim}
\subsection{Verify slurmd}
\label{sec:org89356ca}

\begin{verbatim}
systemctl status slurmd
\end{verbatim}

Expected:

\begin{verbatim}
active (running)
\end{verbatim}
\subsection{Reconfigure Slurm}
\label{sec:org54b5db5}

On master:

\begin{verbatim}
sudo scontrol reconfigure
\end{verbatim}
\subsection{IMPORTANT Issue Encountered — Nodes in DRAIN State}
\label{sec:orgc10a0b6}

Problem:

Nodes entered:

\begin{verbatim}
DRAIN
\end{verbatim}

after CPU topology changes.

Reason:

Slurm detected mismatch between:
\begin{itemize}
\item configured CPUs
\item detected hardware topology
\end{itemize}

Resolution:

Clear stale slurmd state.
\subsection{Stop slurmd}
\label{sec:orgd885c68}

On compute nodes:

\begin{verbatim}
sudo systemctl stop slurmd
\end{verbatim}
\subsection{Remove Old State}
\label{sec:org401d987}

\begin{verbatim}
sudo rm -rf /var/spool/slurmd/*
\end{verbatim}
\subsection{Restart slurmd}
\label{sec:org613c576}

\begin{verbatim}
sudo systemctl start slurmd
\end{verbatim}
\subsection{Resume Nodes}
\label{sec:org385c821}

On master:

\begin{verbatim}
sudo scontrol update NodeName=compute01 State=RESUME
\end{verbatim}

\begin{verbatim}
sudo scontrol update NodeName=compute02 State=RESUME
\end{verbatim}

\begin{verbatim}
sudo scontrol update NodeName=compute03 State=RESUME
\end{verbatim}
\subsection{Verify Cluster Status}
\label{sec:orga3224b5}

\begin{verbatim}
sinfo
\end{verbatim}

Expected:

\begin{verbatim}
PARTITION AVAIL TIMELIMIT NODES STATE NODELIST
debug* up infinite 3 idle compute[01-03]
\end{verbatim}
\subsection{Verify Node Details}
\label{sec:org1e8ba79}

\begin{verbatim}
scontrol show nodes
\end{verbatim}

This displays:

\begin{itemize}
\item CPU information
\item memory
\item node states
\item scheduler topology
\item resource allocation
\end{itemize}
\subsection{Test Slurm Execution}
\label{sec:org6c2879b}
\begin{enumerate}
\item Single Node Test
\label{sec:orgacb0e75}

From login node:

\begin{verbatim}
srun hostname
\end{verbatim}

Expected:

\begin{verbatim}
compute01.local
\end{verbatim}
\end{enumerate}
\subsection{Multi-node Execution}
\label{sec:orgd8fcb22}

\begin{verbatim}
srun -N 3 hostname
\end{verbatim}

Expected:

\begin{verbatim}
compute01.local
compute02.local
compute03.local
\end{verbatim}
\subsection{Multi-task Multi-node Execution}
\label{sec:orga805b25}

\begin{verbatim}
srun -N 3 -n 6 hostname
\end{verbatim}

This verifies:

\begin{itemize}
\item task distribution
\item multi-node scheduling
\item CPU allocation
\end{itemize}
\subsection{Multi-core Single-node Execution}
\label{sec:org08ba8ae}

\begin{verbatim}
srun -N 1 -n 2 hostname
\end{verbatim}

Expected:

\begin{verbatim}
compute01.local
compute01.local
\end{verbatim}
\subsection{IMPORTANT Slurm Scheduling Concept}
\label{sec:orgbaf45c3}

Earlier:

\begin{verbatim}
srun -N 1 -n 2 hostname
\end{verbatim}

failed because nodes initially had:

\begin{verbatim}
CPUs=1
\end{verbatim}

After increasing VM CPUs and updating slurm.conf:

\begin{verbatim}
CPUs=2
\end{verbatim}

the scheduler successfully allocated 2 tasks on a single node.

Important HPC lesson:

Slurm strictly enforces:
\begin{itemize}
\item CPU limits
\item task allocation
\item topology constraints
\item resource consistency
\end{itemize}
\section{Shared Software Filesystem (/apps)}
\label{sec:org74b9ab4}

Goal:

Create a dedicated shared software filesystem for:

\begin{itemize}
\item Spack
\item compilers
\item MPI stacks
\item scientific applications
\item modulefiles
\item HPC software hierarchy
\end{itemize}

This phase transitions the cluster from:

\begin{verbatim}
basic scheduler infrastructure
\end{verbatim}

to:

\begin{verbatim}
realistic HPC software environment
\end{verbatim}
\subsection{Why Separate /apps Filesystem?}
\label{sec:org74ec4d3}

Initially:

\begin{itemize}
\item each VM had a 30 GB local disk
\item /home was exported from master over NFS
\end{itemize}

Concern encountered:

\begin{verbatim}
Shared /home only had ~30 GB available.
Spack and HPC software stacks can become very large.
\end{verbatim}

Important HPC storage concept:

Real HPC clusters usually separate:

\begin{center}
\begin{tabular}{ll}
Filesystem & Purpose\\
\hline
/home & user files\\
/apps & software stack\\
/scratch & temporary runtime data\\
/archive & long-term storage\\
\end{tabular}
\end{center}

Therefore a dedicated:

\begin{verbatim}
/apps
\end{verbatim}

filesystem was created.

This is significantly closer to real HPC architecture.
\subsection{IMPORTANT Storage Architecture Concept}
\label{sec:orgf529c17}

Each VM still has its own independent local storage.

Example:

\begin{center}
\begin{tabular}{ll}
Node & Local Disk\\
\hline
master & 30 GB\\
login & 30 GB\\
compute01 & 30 GB\\
compute02 & 30 GB\\
compute03 & 30 GB\\
\end{tabular}
\end{center}

When:

\begin{verbatim}
master:/home
\end{verbatim}

was mounted via NFS:

\begin{itemize}
\item compute node local /home became hidden
\item but local disks still exist
\end{itemize}

Only the exported filesystem becomes shared.

Therefore:

\begin{verbatim}
/apps
\end{verbatim}

was created as a dedicated shared software layer.
\subsection{Add New Virtual Disk to Master}
\label{sec:org4f2109e}

A second virtual disk was attached ONLY to:

\begin{verbatim}
master
\end{verbatim}

Purpose:

\begin{itemize}
\item centralized software storage
\item shared application filesystem
\item Spack software stack
\end{itemize}

Disk size selected:

\begin{verbatim}
60 GB
\end{verbatim}

Disk type:

\begin{verbatim}
VDI (Dynamically Allocated)
\end{verbatim}
\subsection{Verify New Disk Detection}
\label{sec:org6c568b0}

After booting master:

\begin{verbatim}
lsblk
\end{verbatim}

Expected:

\begin{verbatim}
sda    30G
sdb    60G
\end{verbatim}

Important:

\begin{center}
\begin{tabular}{ll}
Device & Purpose\\
\hline
sda & OS disk\\
sdb & New apps disk\\
\end{tabular}
\end{center}

Only:

\begin{verbatim}
sdb
\end{verbatim}

was modified.
\subsection{Create GPT Partition Table}
\label{sec:orge906923}

\begin{verbatim}
sudo fdisk /dev/sdb
\end{verbatim}

Inside fdisk:

Create GPT label:

\begin{verbatim}
g
\end{verbatim}

Create partition:

\begin{verbatim}
n
\end{verbatim}

Accept defaults.

Write changes:

\begin{verbatim}
w
\end{verbatim}
\subsection{Verify Partition}
\label{sec:org9c1d581}

\begin{verbatim}
lsblk
\end{verbatim}

Expected:

\begin{verbatim}
sdb
└── sdb1
\end{verbatim}
\subsection{Create XFS Filesystem}
\label{sec:orgceb6b76}

\begin{verbatim}
sudo mkfs.xfs /dev/sdb1
\end{verbatim}

Why XFS?

\begin{itemize}
\item default enterprise filesystem
\item highly scalable
\item excellent for large filesystems
\item commonly used in HPC
\end{itemize}
\subsection{Create Mount Point}
\label{sec:org33ae470}

\begin{verbatim}
sudo mkdir -p /apps
\end{verbatim}
\subsection{Get Filesystem UUID}
\label{sec:org1b06050}

\begin{verbatim}
sudo blkid /dev/sdb1
\end{verbatim}

Example:

\begin{verbatim}
UUID="xxxxxxxx-xxxx-xxxx"
\end{verbatim}
\subsection{Configure Persistent Mount}
\label{sec:org792e1ca}

Edit:

\begin{verbatim}
sudo vim /etc/fstab
\end{verbatim}

Add:

\begin{verbatim}
UUID=YOUR_UUID_HERE   /apps   xfs   defaults   0 0
\end{verbatim}

Important Linux administration concept:

UUIDs are preferred over:

\begin{verbatim}
/dev/sdb1
\end{verbatim}

because device names can change after reboot.
\subsection{Mount Filesystem}
\label{sec:orgc02d77f}

\begin{verbatim}
sudo mount -a
\end{verbatim}
\subsection{Verify Mount}
\label{sec:org3e14414}

\begin{verbatim}
df -h
\end{verbatim}

Expected:

\begin{verbatim}
/apps
\end{verbatim}

mounted with approximately:

\begin{verbatim}
60 GB
\end{verbatim}
\subsection{Configure Ownership}
\label{sec:org6aefedf}

\begin{verbatim}
sudo chown hpcadmin:hpcadmin /apps
\end{verbatim}
\subsection{Create Shared HPC Software Hierarchy}
\label{sec:orga0a9e3c}

\begin{verbatim}
mkdir -p /apps/{spack,modulefiles,software,sources}
\end{verbatim}

Resulting structure:

\begin{verbatim}
/apps
├── spack
├── modulefiles
├── software
└── sources
\end{verbatim}
\subsection{HPC Software Hierarchy Concept}
\label{sec:orgd868d9f}

Purpose of directories:

\begin{center}
\begin{tabular}{ll}
Directory & Purpose\\
\hline
spack & Spack package manager\\
modulefiles & custom modulefiles\\
software & installed software\\
sources & downloaded sources\\
\end{tabular}
\end{center}

This is similar to production HPC software layouts.
\subsection{Export /apps via NFS}
\label{sec:orge6c9d30}

Edit exports file:

\begin{verbatim}
sudo vim /etc/exports
\end{verbatim}

Add:

\begin{verbatim}
/apps 10.10.10.0/24(rw,sync,no_root_squash)
\end{verbatim}
\subsection{Apply NFS Exports}
\label{sec:org7257f01}

\begin{verbatim}
sudo exportfs -rav
\end{verbatim}
\subsection{Verify Exports}
\label{sec:orgb01a3ae}

\begin{verbatim}
showmount -e localhost
\end{verbatim}

Expected:

\begin{verbatim}
/apps
/home
\end{verbatim}
\subsection{Create Mount Points on Client Nodes}
\label{sec:org4f048fd}

On:
\begin{itemize}
\item login
\item compute01
\item compute02
\item compute03
\end{itemize}

Create:

\begin{verbatim}
sudo mkdir -p /apps
\end{verbatim}
\subsection{Configure Persistent NFS Mounts}
\label{sec:org083296f}

Edit:

\begin{verbatim}
sudo vim /etc/fstab
\end{verbatim}

Add:

\begin{verbatim}
master:/apps   /apps   nfs   defaults,_netdev   0 0
\end{verbatim}
\subsection{Mount Shared /apps Filesystem}
\label{sec:org3f81716}

\begin{verbatim}
sudo mount -a
\end{verbatim}
\subsection{Verify Shared Mount}
\label{sec:org4f552cb}

\begin{verbatim}
df -h
\end{verbatim}

Expected:

\begin{verbatim}
master:/apps
\end{verbatim}

mounted on all nodes.
\subsection{Verify Shared Filesystem Functionality}
\label{sec:org7159544}

On master:

\begin{verbatim}
touch /apps/testfile
\end{verbatim}

On compute node:

\begin{verbatim}
ls /apps
\end{verbatim}

Expected:

\begin{verbatim}
testfile
\end{verbatim}

This verifies:

\begin{itemize}
\item shared software filesystem
\item NFS functionality
\item cluster-wide software visibility
\end{itemize}
\subsection{HPC Architecture After /apps Setup}
\label{sec:org7bef2d5}

Current cluster storage architecture:

\begin{center}
\begin{tabular}{ll}
Filesystem & Purpose\\
\hline
/home & user files\\
/apps & software stack\\
local disks & local runtime storage\\
\end{tabular}
\end{center}

This resembles real HPC environments significantly more closely.
\section{Spack Software Stack + MPI Runtime Validation}
\label{sec:org3157054}

Goal:

Create centralized HPC software stack using:

\begin{itemize}
\item Spack
\item shared software filesystem
\item shared compilers
\item shared MPI stack
\item scheduler-aware runtime execution
\end{itemize}

This phase validates:

\begin{itemize}
\item software distribution
\item multinode execution
\item Slurm integration
\item MPI runtime functionality
\end{itemize}

After completing this section you will be able to:

\begin{itemize}
\item install HPC software using Spack
\item manage multiple compiler versions
\item manage multiple MPI versions
\item load software dynamically
\item compile MPI applications
\item execute distributed jobs using Slurm
\item monitor and debug jobs
\item understand performance metrics
\end{itemize}
\subsection{IMPORTANT HPC Software Stack Concept}
\label{sec:org32d5249}

Instead of:

\begin{verbatim}
dnf install openmpi
\end{verbatim}

a centralized HPC software stack was created.

Benefits:

\begin{itemize}
\item software installed once
\item accessible cluster-wide
\item compiler hierarchy support
\item reproducible builds
\item module-based software environments
\item centralized software management
\end{itemize}

This approach is significantly closer to real HPC centers.

In production HPC environments:

\begin{itemize}
\item operating system packages are separated
\item scientific software stack is separated
\item compiler stacks are separated
\item MPI stacks are separated
\end{itemize}

This separation improves:

\begin{itemize}
\item portability
\item reproducibility
\item scalability
\item maintenance
\end{itemize}
\subsection{IMPORTANT Shared Filesystem Concept}
\label{sec:org900f4ee}

The cluster uses:

\begin{verbatim}
/apps
\end{verbatim}

as shared application storage.

Because:

\begin{verbatim}
/apps
\end{verbatim}

is mounted cluster-wide through NFS:

\begin{itemize}
\item software installed once becomes visible everywhere
\item compute nodes automatically access installed software
\item users share common software stack
\item recompilation is unnecessary on compute nodes
\end{itemize}

This is one of the MOST important HPC architecture concepts.
\subsection{Install Development Tools}
\label{sec:orgcdce2fa}

Installed manually on ALL nodes:

\begin{verbatim}
sudo dnf groupinstall -y "Development Tools"
\end{verbatim}

Purpose:

Provides:

\begin{itemize}
\item gcc
\item g++
\item gfortran
\item make
\item linker tools
\item build utilities
\item development environment
\end{itemize}

Important note:

\begin{verbatim}
System compilers are still required even when using Spack.
\end{verbatim}

Reason:

Spack builds software from source code and therefore still requires:

\begin{itemize}
\item assembler
\item linker
\item basic compiler toolchain
\end{itemize}
\subsection{Move into Shared Application Filesystem}
\label{sec:orgfe241f8}

Run on:

\begin{verbatim}
master
\end{verbatim}

Move into shared applications directory:

\begin{verbatim}
cd /apps
\end{verbatim}

Verify:

\begin{verbatim}
pwd
\end{verbatim}

Expected:

\begin{verbatim}
/apps
\end{verbatim}
\subsection{Clone Spack into Shared Filesystem}
\label{sec:org542515e}

Clone Spack:

\begin{verbatim}
git clone https://github.com/spack/spack.git
\end{verbatim}

Expected:

\begin{verbatim}
/apps/spack
\end{verbatim}

This creates centralized package manager installation.
\subsection{Enter Spack Directory}
\label{sec:orgf656a92}

\begin{verbatim}
cd spack
\end{verbatim}

Verify:

\begin{verbatim}
pwd
\end{verbatim}

Expected:

\begin{verbatim}
/apps/spack
\end{verbatim}
\subsection{Initialize Spack Environment}
\label{sec:orgb218b6b}

Initialize Spack:

\begin{verbatim}
source share/spack/setup-env.sh
\end{verbatim}

Purpose:

Adds:

\begin{itemize}
\item spack command
\item shell integration
\item environment variables
\item package management environment
\end{itemize}

to current shell session.

Without initialization:

\begin{itemize}
\item spack command will not work
\item spack load will fail
\item software environments will not activate
\end{itemize}
\subsection{IMPORTANT Shell Environment Concept}
\label{sec:orgc4579ef}

\begin{verbatim}
source
\end{verbatim}

does NOT execute script normally.

Instead:

\begin{verbatim}
it modifies current shell environment
\end{verbatim}

This is extremely important in:

\begin{itemize}
\item Spack
\item modules
\item MPI environments
\item compiler environments
\end{itemize}
\subsection{Make Spack Permanent}
\label{sec:org7177105}

To avoid sourcing manually every session:

Edit:

\begin{verbatim}
vim ~/.bashrc
\end{verbatim}

Add:

\begin{verbatim}
source /apps/spack/share/spack/setup-env.sh
\end{verbatim}

Reload shell:

\begin{verbatim}
source ~/.bashrc
\end{verbatim}

Verify:

\begin{verbatim}
spack --version
\end{verbatim}

Expected:
\begin{itemize}
\item Spack version information
\end{itemize}
\subsection{Clone Additional Spack Repository}
\label{sec:org2af2a00}

Clone repository:

\begin{verbatim}
git clone https://github.com/spack/spack-packages.git
\end{verbatim}

Purpose:

Adds additional package definitions.
\subsection{Configure Spack Repository}
\label{sec:org8958134}

\begin{verbatim}
spack repo set --destination "$(pwd)/spack-packages" builtin
\end{verbatim}

Purpose:

Configure additional repository support.

This allows:

\begin{itemize}
\item package metadata
\item package recipes
\item software definitions
\end{itemize}

to become visible to Spack.
\subsection{Verify Spack Functionality}
\label{sec:org8515436}

\begin{verbatim}
spack info gcc
\end{verbatim}

Expected:

Detailed GCC package information.

This verifies:

\begin{itemize}
\item Spack environment working
\item package repository functional
\item package metadata accessible
\end{itemize}
\subsection{IMPORTANT HPC Software Stack Concept}
\label{sec:org736ef3a}

Spack operates as:

\begin{verbatim}
source-based HPC package manager
\end{verbatim}

Unlike:

\begin{center}
\begin{tabular}{ll}
Package Manager & Purpose\\
\hline
dnf & operating system\\
Spack & HPC software ecosystem\\
\end{tabular}
\end{center}

Production HPC systems commonly separate:

\begin{itemize}
\item OS packages
\item scientific software stacks
\end{itemize}

exactly like this setup.
\subsection{IMPORTANT Spack Build Concept}
\label{sec:orgf398b36}

When installing software using Spack:

\begin{verbatim}
software is compiled from source
\end{verbatim}

Advantages:

\begin{itemize}
\item architecture optimization
\item compiler flexibility
\item dependency control
\item reproducibility
\item variant support
\end{itemize}

Disadvantages:

\begin{itemize}
\item longer installation time
\end{itemize}
\subsection{Detect Available Compilers}
\label{sec:orge885e23}

Before installing software:

\begin{verbatim}
spack compiler find
\end{verbatim}

Purpose:

Detect system compilers available for building software.

Verify detected compilers:

\begin{verbatim}
spack compilers
\end{verbatim}

Expected:
\begin{itemize}
\item compiler list
\end{itemize}
\subsection{Install GCC Using Spack}
\label{sec:orgfb7b505}

Install GCC version:

\begin{verbatim}
spack install gcc@13.4.0
\end{verbatim}

Purpose:

Create dedicated HPC compiler stack.

This may take considerable time because:

\begin{itemize}
\item GCC itself is compiled from source
\item dependencies are built
\item optimization occurs
\end{itemize}
\subsection{IMPORTANT Multiple Compiler Version Concept}
\label{sec:orgf990eb4}

Production HPC systems frequently maintain multiple compiler versions simultaneously.

Example:

\begin{verbatim}
gcc@11
gcc@12
gcc@13
intel
nvhpc
clang
\end{verbatim}

Reasons:

\begin{itemize}
\item application compatibility
\item compiler ABI differences
\item reproducibility
\item benchmark consistency
\end{itemize}
\subsection{Verify Installed Software}
\label{sec:orgb325ecf}

List installed packages:

\begin{verbatim}
spack find
\end{verbatim}

Expected example:

\begin{verbatim}
==> 1 installed package
gcc@13.4.0
\end{verbatim}
\subsection{Load Installed GCC Compiler}
\label{sec:org1c14b36}

\begin{verbatim}
spack load gcc@13.4.0
\end{verbatim}

Purpose:

Activate compiler environment.

Verify compiler path:

\begin{verbatim}
which gcc
\end{verbatim}

Expected:
\begin{itemize}
\item Spack GCC path
\end{itemize}

Verify version:

\begin{verbatim}
gcc --version
\end{verbatim}

Expected:
\begin{itemize}
\item GCC 13.4.0
\end{itemize}
\subsection{IMPORTANT Runtime Environment Concept}
\label{sec:org7eb4a66}

When software is loaded:

\begin{itemize}
\item PATH changes
\item LD\_LIBRARY\_PATH changes
\item MANPATH changes
\end{itemize}

This allows applications to locate:

\begin{itemize}
\item executables
\item shared libraries
\item runtime dependencies
\end{itemize}

automatically.
\subsection{Register Newly Installed Compiler}
\label{sec:org17554ff}

After installing GCC:

\begin{verbatim}
spack compiler find
\end{verbatim}

Purpose:

Allow Spack to use newly installed GCC for future package builds.

Verify:

\begin{verbatim}
spack compilers
\end{verbatim}

Expected:
\begin{itemize}
\item new GCC compiler entry
\end{itemize}
\subsection{Install OpenMPI Using Spack}
\label{sec:orgc2a76b3}

Install MPI runtime:

\begin{verbatim}
spack install openmpi
\end{verbatim}

Purpose:

Create shared MPI runtime stack.

This installation includes:

\begin{itemize}
\item MPI runtime
\item MPI compiler wrappers
\item MPI libraries
\item distributed runtime environment
\end{itemize}
\subsection{IMPORTANT MPI Compiler Dependency Concept}
\label{sec:orgf45a0b8}

MPI stacks are tightly coupled with compiler toolchains.

Example:

\begin{verbatim}
openmpi built using gcc@13
\end{verbatim}

This is important because:

\begin{itemize}
\item ABI compatibility matters
\item runtime compatibility matters
\item mixed compiler environments may fail
\end{itemize}
\subsection{Verify Installed Software Tree}
\label{sec:orga712175}

Show dependency tree:

\begin{verbatim}
spack find -d
\end{verbatim}

Purpose:

Displays:

\begin{itemize}
\item dependency hierarchy
\item compiler hierarchy
\item runtime dependencies
\end{itemize}

Example:

\begin{verbatim}
openmpi
   ^gcc@13.4.0
\end{verbatim}
\subsection{IMPORTANT Spack Spec Concept}
\label{sec:org805f40f}

Every Spack installation contains:

\begin{itemize}
\item package version
\item compiler information
\item dependency tree
\item unique installation hash
\end{itemize}

Example:

\begin{verbatim}
openmpi@5.0.7%gcc@13.4.0
\end{verbatim}

Meaning:

\begin{center}
\begin{tabular}{ll}
Component & Meaning\\
\hline
openmpi & package\\
5.0.7 & package version\\
gcc@13.4.0 & compiler used\\
\end{tabular}
\end{center}
\subsection{IMPORTANT Hash-Based Software Loading}
\label{sec:org1e77e49}

Sometimes multiple builds of same version exist because of:

\begin{itemize}
\item different compilers
\item different variants
\item different dependencies
\end{itemize}

Spack assigns:

\begin{verbatim}
unique hashes
\end{verbatim}

View hashes:

\begin{verbatim}
spack find -l
\end{verbatim}

Expected:

\begin{verbatim}
abc1234 gcc@13.4.0
xyz5678 openmpi@5.0.7
\end{verbatim}

Load using hash:

\begin{verbatim}
spack load /xyz5678
\end{verbatim}

This is extremely common in production HPC systems.
\subsection{Load OpenMPI Environment}
\label{sec:org0d322db}

\begin{verbatim}
spack load openmpi
\end{verbatim}

Verify:

\begin{verbatim}
which mpicc
\end{verbatim}

Expected:
\begin{itemize}
\item Spack OpenMPI wrapper path
\end{itemize}

Verify MPI runtime:

\begin{verbatim}
mpirun --version
\end{verbatim}

Expected:
\begin{itemize}
\item OpenMPI version information
\end{itemize}
\subsection{IMPORTANT MPI Wrapper Concept}
\label{sec:org7310109}

\begin{verbatim}
mpicc
\end{verbatim}

is NOT a compiler itself.

It is a wrapper that automatically injects:

\begin{itemize}
\item MPI headers
\item MPI libraries
\item runtime linker flags
\end{itemize}

around GCC.

This is a critical HPC build concept.
\subsection{Verify Loaded Software}
\label{sec:org59ec366}

Display loaded packages:

\begin{verbatim}
spack find --loaded
\end{verbatim}

Unload OpenMPI:

\begin{verbatim}
spack unload openmpi
\end{verbatim}

Unload all software:

\begin{verbatim}
spack unload --all
\end{verbatim}

Verify:

\begin{verbatim}
spack find --loaded
\end{verbatim}

Expected:
\begin{itemize}
\item no loaded packages
\end{itemize}
\subsection{Create MPI Performance Validation Program}
\label{sec:org150fdba}

Create source file:

\begin{verbatim}
vim mpi_sum.c
\end{verbatim}

Purpose:

\begin{itemize}
\item validate MPI runtime
\item validate multinode execution
\item validate scheduler-aware execution
\item compare serial vs parallel runtime
\item demonstrate reduction operations
\item calculate speedup
\item calculate efficiency
\end{itemize}

Paste:

\begin{verbatim}
#include <mpi.h>
#include <stdio.h>
#include <stdlib.h>

#define N 1000000000

double serial_sum()
{
    double sum = 0.0;

    for(long long int i = 0; i < N; i++)
    {
        sum += i;
    }

    return sum;
}

int main(int argc, char *argv[])
{
    int rank, size;

    long long int start_index;
    long long int end_index;
    long long int chunk;

    double local_sum = 0.0;
    double global_sum = 0.0;

    double serial_start;
    double serial_end;

    double parallel_start;
    double parallel_end;

    MPI_Init(&argc, &argv);

    MPI_Comm_rank(MPI_COMM_WORLD, &rank);
    MPI_Comm_size(MPI_COMM_WORLD, &size);

    if(rank == 0)
    {
        serial_start = MPI_Wtime();

        double serial_result = serial_sum();

        serial_end = MPI_Wtime();

        printf("\n");
        printf("========== SERIAL COMPUTATION ==========\n");
        printf("Sum = %.0f\n", serial_result);
        printf("Time taken = %f seconds\n",
               serial_end - serial_start);
    }

    MPI_Barrier(MPI_COMM_WORLD);

    chunk = N / size;

    start_index = rank * chunk;

    if(rank == size - 1)
        end_index = N;
    else
        end_index = start_index + chunk;

    parallel_start = MPI_Wtime();

    for(long long int i = start_index;
        i < end_index;
        i++)
    {
        local_sum += i;
    }

    MPI_Reduce(&local_sum,
               &global_sum,
               1,
               MPI_DOUBLE,
               MPI_SUM,
               0,
               MPI_COMM_WORLD);

    parallel_end = MPI_Wtime();

    if(rank == 0)
    {
        double parallel_time =
            parallel_end - parallel_start;

        double serial_time =
            serial_end - serial_start;

        printf("\n");
        printf("========== PARALLEL COMPUTATION ==========\n");
        printf("Processes used = %d\n", size);
        printf("Sum = %.0f\n", global_sum);
        printf("Time taken = %f seconds\n",
               parallel_time);

        printf("\n");
        printf("========== PERFORMANCE ==========\n");
        printf("Speedup = %f\n",
               serial_time / parallel_time);

        printf("Efficiency = %f\n",
               (serial_time / parallel_time) / size);
    }

    MPI_Finalize();

    return 0;
}
\end{verbatim}
\subsection{IMPORTANT MPI Reduction Concept}
\label{sec:orgef9a491}

This program uses:

\begin{verbatim}
MPI_Reduce()
\end{verbatim}

Purpose:

Combine partial sums from all MPI ranks into:

\begin{verbatim}
single global result
\end{verbatim}

Flow:

\begin{verbatim}
local sums
    ↓
MPI_Reduce
    ↓
global sum
\end{verbatim}

This is one of the MOST important MPI collective communication operations.
\subsection{IMPORTANT Parallel Execution Concept}
\label{sec:orgfa1aca0}

Serial execution:

\begin{verbatim}
single process computes entire workload
\end{verbatim}

Parallel execution:

\begin{verbatim}
multiple MPI ranks divide workload
\end{verbatim}

Example:

\begin{center}
\begin{tabular}{ll}
Rank & Work\\
\hline
rank0 & 0 --> 25\%\\
rank1 & 25\% --> 50\%\\
rank2 & 50\% --> 75\%\\
rank3 & 75\% --> 100\%\\
\end{tabular}
\end{center}

Advantages:

\begin{itemize}
\item reduced execution time
\item distributed workload
\item scalable computation
\end{itemize}
\subsection{Load Software Before Compilation}
\label{sec:org089743e}

Initialize Spack:

\begin{verbatim}
source /apps/spack/share/spack/setup-env.sh
\end{verbatim}

Load compiler:

\begin{verbatim}
spack load gcc@13.4.0
\end{verbatim}

Load MPI runtime:

\begin{verbatim}
spack load openmpi
\end{verbatim}

Verify:

\begin{verbatim}
which mpicc
\end{verbatim}

Expected:
\begin{itemize}
\item MPI compiler wrapper path
\end{itemize}
\subsection{Compile MPI Program}
\label{sec:orgd5c0b60}

Compile:

\begin{verbatim}
mpicc mpi_sum.c -o mpi_sum
\end{verbatim}

Verify executable:

\begin{verbatim}
ls
\end{verbatim}

Expected:

\begin{verbatim}
mpi_sum
\end{verbatim}
\subsection{IMPORTANT MPI Compilation Concept}
\label{sec:org9e0c8ed}

MPI programs are compiled using:

\begin{verbatim}
mpicc
\end{verbatim}

instead of:

\begin{verbatim}
gcc
\end{verbatim}

because MPI wrapper automatically injects:

\begin{itemize}
\item MPI headers
\item MPI libraries
\item linker flags
\end{itemize}
\subsection{Interactive MPI Runtime Validation}
\label{sec:orgffb3308}

Run interactively:

\begin{verbatim}
mpirun -np 4 ./mpi_sum
\end{verbatim}

Expected:

\begin{verbatim}
========== SERIAL COMPUTATION ==========
...

========== PARALLEL COMPUTATION ==========
...

========== PERFORMANCE ==========
...
\end{verbatim}

This validates:

\begin{itemize}
\item MPI runtime
\item MPI communication
\item reduction operations
\item distributed execution
\end{itemize}
\subsection{IMPORTANT Shared Filesystem Advantage}
\label{sec:orgb2af4a5}

Because:

\begin{itemize}
\item source code
\item executable
\item software stack
\end{itemize}

exist on shared filesystem:

\begin{verbatim}
/apps
or
/home
\end{verbatim}

compute nodes automatically access:

\begin{itemize}
\item compiled binaries
\item runtime libraries
\item MPI stack
\end{itemize}

without recompilation.

This is one of the major HPC architecture advantages.
\subsection{Create Slurm Batch Script}
\label{sec:org2a66b9c}

Create batch script:

\begin{verbatim}
vim submit.sh
\end{verbatim}

Paste:

\begin{verbatim}
#!/bin/bash
#SBATCH -N 3
#SBATCH --ntasks-per-node=1
#SBATCH --job-name=test
#SBATCH --output=%J.out
#SBATCH --error=%J.err
#SBATCH --time=01:00:00
##SBATCH --nodelist=compute03,compute01

source /apps/spack/share/spack/setup-env.sh

spack load openmpi

mpirun -np $SLURM_NTASKS ./$1
\end{verbatim}

Make executable:

\begin{verbatim}
chmod +x submit.sh
\end{verbatim}
\subsection{IMPORTANT Slurm Script Explanation}
\label{sec:orgf5d275b}

\begin{center}
\begin{tabular}{ll}
Line & Purpose\\
\hline
\#SBATCH -N 3 & request 3 nodes\\
--ntasks-per-node=1 & one MPI rank per node\\
--job-name=test & job name\\
--output=\%J.out & stdout file\\
--error=\%J.err & stderr file\\
--time=01:00:00 & walltime limit\\
\end{tabular}
\end{center}
\subsection{IMPORTANT Runtime Environment Section}
\label{sec:org3a642d0}

Inside batch script:

\begin{verbatim}
source /apps/spack/share/spack/setup-env.sh
\end{verbatim}

initializes Spack environment.

\begin{verbatim}
spack load openmpi
\end{verbatim}

loads MPI runtime stack.

Without these:

\begin{itemize}
\item mpirun may fail
\item MPI libraries may not load
\item runtime linker errors may occur
\end{itemize}
\subsection{IMPORTANT Slurm Runtime Variable}
\label{sec:orgaf68072}

\begin{verbatim}
$SLURM_NTASKS
\end{verbatim}

contains:

\begin{verbatim}
total MPI process count allocated by Slurm
\end{verbatim}

Example:

\begin{center}
\begin{tabular}{rrr}
Nodes & Tasks per node & Total\\
\hline
3 & 1 & 3\\
\end{tabular}
\end{center}

Therefore:

\begin{verbatim}
mpirun -np $SLURM_NTASKS
\end{verbatim}

automatically launches correct number of MPI ranks.
\subsection{Submit MPI Job}
\label{sec:org2ff15c5}

Submit job:

\begin{verbatim}
sbatch submit.sh mpi_sum
\end{verbatim}

Expected:

\begin{verbatim}
Submitted batch job 101
\end{verbatim}

This validates:

\begin{itemize}
\item Slurm scheduler
\item distributed allocation
\item multinode MPI execution
\item scheduler-aware runtime launch
\end{itemize}
\subsection{Monitor Jobs}
\label{sec:orge69c20c}

Check user jobs:

\begin{verbatim}
squeue --me
\end{verbatim}

Expected:

\begin{verbatim}
JOBID PARTITION NAME USER ST TIME NODES NODELIST(REASON)
\end{verbatim}
\subsection{IMPORTANT Job States}
\label{sec:org0799f86}

\begin{center}
\begin{tabular}{ll}
State & Meaning\\
\hline
PD & pending\\
R & running\\
CG & completing\\
CD & completed\\
F & failed\\
\end{tabular}
\end{center}
\subsection{View Detailed Job Information}
\label{sec:org2cc9940}

Show detailed job info:

\begin{verbatim}
scontrol show job <jobid>
\end{verbatim}

Example:

\begin{verbatim}
scontrol show job 101
\end{verbatim}

Displays:

\begin{itemize}
\item allocated nodes
\item CPUs
\item runtime
\item scheduling information
\item execution directory
\item output locations
\end{itemize}
\subsection{Check Cluster Status}
\label{sec:org1db87d5}

Check cluster nodes:

\begin{verbatim}
sinfo
\end{verbatim}

Expected:

\begin{verbatim}
debug* up infinite 3 idle compute[01-03]
\end{verbatim}
\subsection{IMPORTANT Node States}
\label{sec:orgdcf24f1}

\begin{center}
\begin{tabular}{ll}
State & Meaning\\
\hline
idle & available\\
alloc & allocated\\
mix & partially used\\
drain & unavailable\\
down & offline\\
\end{tabular}
\end{center}
\subsection{Check Output Files}
\label{sec:orgf264212}

After completion:

\begin{verbatim}
ls
\end{verbatim}

Expected:

\begin{verbatim}
101.out
101.err
\end{verbatim}

View output:

\begin{verbatim}
cat 101.out
\end{verbatim}

Expected:

\begin{itemize}
\item serial runtime
\item parallel runtime
\item speedup
\item efficiency
\end{itemize}
\subsection{IMPORTANT HPC Performance Metrics}
\label{sec:org8ad284f}

The program calculates:

\begin{center}
\begin{tabular}{ll}
Metric & Meaning\\
\hline
Speedup & serial\_time / parallel\_time\\
Efficiency & speedup / processes\\
\end{tabular}
\end{center}

Ideal values:

\begin{center}
\begin{tabular}{ll}
Metric & Ideal\\
\hline
Speedup & near process count\\
Efficiency & near 1.0\\
\end{tabular}
\end{center}

Real systems experience:

\begin{itemize}
\item communication overhead
\item synchronization cost
\item MPI latency
\item reduction overhead
\end{itemize}
\subsection{Cancel Jobs}
\label{sec:org61e0d03}

Cancel job:

\begin{verbatim}
scancel <jobid>
\end{verbatim}

Example:

\begin{verbatim}
scancel 101
\end{verbatim}

Cancel all user jobs:

\begin{verbatim}
scancel -u $USER
\end{verbatim}
\subsection{Interactive Slurm Execution}
\label{sec:org8850fbf}

Launch interactive shell:

\begin{verbatim}
srun --pty bash
\end{verbatim}

Distributed hostname test:

\begin{verbatim}
srun -N 3 -n 3 hostname
\end{verbatim}

Expected:
\begin{itemize}
\item hostnames from allocated nodes
\end{itemize}

This validates:

\begin{itemize}
\item scheduler allocation
\item distributed execution
\item runtime communication
\end{itemize}
\subsection{IMPORTANT HPC Workflow}
\label{sec:orga4a9449}

Typical HPC workflow:

\begin{verbatim}
Load software
    ↓
Compile application
    ↓
Prepare Slurm script
    ↓
Submit job
    ↓
Monitor execution
    ↓
Analyze outputs
\end{verbatim}
\subsection{IMPORTANT Production HPC Concept}
\label{sec:org0c2e2e5}

Production HPC systems usually separate:

\begin{itemize}
\item operating system packages
\item scientific software stack
\item compiler stack
\item MPI runtime stack
\item scheduler integration
\end{itemize}

exactly like this setup.

This architecture provides:

\begin{itemize}
\item reproducibility
\item scalability
\item centralized administration
\item software consistency
\item runtime portability
\end{itemize}
\section{Centralized Identity Management using OpenLDAP + nslcd}
\label{sec:org52efbae}

Goal:

Transform the cluster from:

\begin{verbatim}
local-user-only infrastructure
\end{verbatim}

into:

\begin{verbatim}
centralized multi-user HPC environment
\end{verbatim}

using:

\begin{itemize}
\item OpenLDAP
\item nslcd
\item NSS
\item PAM
\end{itemize}

This phase introduces centralized user management for the HPC cluster.

After completion:

\begin{itemize}
\item users are created once centrally
\item login node recognizes LDAP users
\item compute nodes recognize LDAP users
\item centralized authentication works cluster-wide
\item automatic home directory creation works
\item Slurm can later use centralized identities
\end{itemize}
\subsection{Important Architecture}
\label{sec:org55bc6b2}

Cluster roles:

\begin{center}
\begin{tabular}{ll}
Node & Role\\
\hline
master & LDAP server\\
login & LDAP client\\
compute01 & LDAP client\\
compute02 & LDAP client\\
compute03 & LDAP client\\
\end{tabular}
\end{center}

Important concept:

\begin{verbatim}
Only master runs LDAP SERVER (slapd).

Login and compute nodes run LDAP CLIENTS (nslcd).
\end{verbatim}
\subsection{Important Linux Authentication Concepts}
\label{sec:orgeca0f29}

Before starting, understand the layers.

\begin{center}
\begin{tabular}{ll}
Component & Purpose\\
\hline
LDAP & centralized user database\\
NSS & user/group lookup\\
PAM & authentication\\
nslcd & bridge between Linux and LDAP\\
\end{tabular}
\end{center}

Authentication flow:

\begin{verbatim}
SSH Login
   ↓
PAM
   ↓
nslcd
   ↓
LDAP server
\end{verbatim}
\subsection{Install OpenLDAP Server on Master Node}
\label{sec:org1124afd}

Run ONLY on:

\begin{verbatim}
master
\end{verbatim}

Install LDAP packages:

\begin{verbatim}
sudo dnf install -y \
openldap \
openldap-servers \
openldap-clients
\end{verbatim}

Package explanation:

\begin{center}
\begin{tabular}{ll}
Package & Purpose\\
\hline
openldap & LDAP utilities\\
openldap-servers & slapd server daemon\\
openldap-clients & ldapsearch, ldapadd etc\\
\end{tabular}
\end{center}
\subsection{Enable LDAP Server}
\label{sec:org06a3927}

Run ONLY on:

\begin{verbatim}
master
\end{verbatim}

Enable and start slapd:

\begin{verbatim}
sudo systemctl enable --now slapd
\end{verbatim}

Verify:

\begin{verbatim}
systemctl status slapd
\end{verbatim}

Expected:

\begin{verbatim}
active (running)
\end{verbatim}
\subsection{Generate LDAP Admin Password Hash}
\label{sec:orgf451a89}

Run ONLY on:

\begin{verbatim}
master
\end{verbatim}

Generate secure LDAP admin password hash:

\begin{verbatim}
slappasswd
\end{verbatim}

Enter desired password.

Example:

\begin{verbatim}
cluster123
\end{verbatim}

Example output:

\begin{verbatim}
{SSHA}xxxxxxxxxxxxxxxx
\end{verbatim}

Important:

\begin{verbatim}
Copy this hash carefully.
\end{verbatim}

It will be used in LDAP database configuration.
\subsection{Configure LDAP Database}
\label{sec:org1270353}

Run ONLY on:

\begin{verbatim}
master
\end{verbatim}

Create LDAP database configuration file:

\begin{verbatim}
vim ~/db.ldif
\end{verbatim}

Paste:

\begin{verbatim}
dn: olcDatabase={2}mdb,cn=config
changetype: modify
replace: olcSuffix
olcSuffix: dc=hpc,dc=local

dn: olcDatabase={2}mdb,cn=config
changetype: modify
replace: olcRootDN
olcRootDN: cn=Manager,dc=hpc,dc=local

dn: olcDatabase={2}mdb,cn=config
changetype: modify
replace: olcRootPW
olcRootPW: PASTE_HASH_HERE
\end{verbatim}

Replace:

\begin{verbatim}
PASTE_HASH_HERE
\end{verbatim}

with the hash generated using:

\begin{verbatim}
slappasswd
\end{verbatim}
\subsection{Important LDAP Concepts}
\label{sec:org709fe85}

LDAP uses a hierarchical tree called:

\begin{verbatim}
DIT -- Directory Information Tree
\end{verbatim}

Example structure:

\begin{verbatim}
dc=hpc,dc=local
 ├── ou=People
 └── ou=Groups
\end{verbatim}

Meaning of terms:

\begin{center}
\begin{tabular}{ll}
Term & Meaning\\
\hline
dc & domain component\\
ou & organizational unit\\
cn & common name\\
\end{tabular}
\end{center}
\subsection{Apply LDAP Database Configuration}
\label{sec:orga168289}

Run ONLY on:

\begin{verbatim}
master
\end{verbatim}

Apply configuration:

\begin{verbatim}
sudo ldapmodify -Y EXTERNAL -H ldapi:/// \
-f ~/db.ldif
\end{verbatim}

Expected:

\begin{verbatim}
modifying entry
\end{verbatim}
\subsection{Import LDAP Schemas}
\label{sec:org9fce33b}

Schemas define:

\begin{itemize}
\item user attributes
\item group attributes
\item UNIX account structure
\end{itemize}

Run ONLY on:

\begin{verbatim}
master
\end{verbatim}

Import cosine schema:

\begin{verbatim}
sudo ldapadd -Y EXTERNAL -H ldapi:/// \
-f /etc/openldap/schema/cosine.ldif
\end{verbatim}

Import nis schema:

\begin{verbatim}
sudo ldapadd -Y EXTERNAL -H ldapi:/// \
-f /etc/openldap/schema/nis.ldif
\end{verbatim}

Import inetorgperson schema:

\begin{verbatim}
sudo ldapadd -Y EXTERNAL -H ldapi:/// \
-f /etc/openldap/schema/inetorgperson.ldif
\end{verbatim}

Schema explanation:

\begin{center}
\begin{tabular}{ll}
Schema & Purpose\\
\hline
cosine & basic directory attributes\\
nis & UNIX/Linux account support\\
inetorgperson & user/person identity support\\
\end{tabular}
\end{center}
\subsection{Create LDAP Base Tree}
\label{sec:orgc072a35}

Run ONLY on:

\begin{verbatim}
master
\end{verbatim}

Create base LDAP tree:

\begin{verbatim}
vim ~/base.ldif
\end{verbatim}

Paste:

\begin{verbatim}
dn: dc=hpc,dc=local
objectClass: top
objectClass: dcObject
objectClass: organization
o: HPC Cluster
dc: hpc

dn: ou=People,dc=hpc,dc=local
objectClass: organizationalUnit
ou: People

dn: ou=Groups,dc=hpc,dc=local
objectClass: organizationalUnit
ou: Groups
\end{verbatim}

Apply tree:

\begin{verbatim}
ldapadd -x \
-D "cn=Manager,dc=hpc,dc=local" \
-W \
-f ~/base.ldif
\end{verbatim}

Enter LDAP admin password.
\subsection{Verify LDAP Tree}
\label{sec:org3721a94}

Run ONLY on:

\begin{verbatim}
master
\end{verbatim}

Verify tree:

\begin{verbatim}
ldapsearch -x -b "dc=hpc,dc=local"
\end{verbatim}

Expected:

\begin{itemize}
\item dc=hpc,dc=local
\item ou=People
\item ou=Groups
\end{itemize}
\subsection{Create LDAP Group}
\label{sec:orge73e9a7}

Run ONLY on:

\begin{verbatim}
master
\end{verbatim}

Create group file:

\begin{verbatim}
vim ~/group.ldif
\end{verbatim}

Paste:

\begin{verbatim}
dn: cn=hpcusers,ou=Groups,dc=hpc,dc=local
objectClass: posixGroup
cn: hpcusers
gidNumber: 10000
\end{verbatim}

Add group:

\begin{verbatim}
ldapadd -x \
-D "cn=Manager,dc=hpc,dc=local" \
-W \
-f ~/group.ldif
\end{verbatim}
\subsection{Create LDAP User Password Hash}
\label{sec:org73cf3e5}

Run ONLY on:

\begin{verbatim}
master
\end{verbatim}

Generate password hash:

\begin{verbatim}
slappasswd
\end{verbatim}

Example password:

\begin{verbatim}
test123
\end{verbatim}

Copy generated hash carefully.
\subsection{Create LDAP User}
\label{sec:org756700e}

Run ONLY on:

\begin{verbatim}
master
\end{verbatim}

Create user file:

\begin{verbatim}
vim ~/user.ldif
\end{verbatim}

Paste:

\begin{verbatim}
dn: uid=testUser,ou=People,dc=hpc,dc=local
objectClass: inetOrgPerson
objectClass: posixAccount
objectClass: shadowAccount
cn: Test User
sn: User
uid: testUser
uidNumber: 10001
gidNumber: 10000
homeDirectory: /home/testUser
loginShell: /bin/bash
userPassword: PASTE_HASH_HERE
\end{verbatim}

Replace:

\begin{verbatim}
PASTE_HASH_HERE
\end{verbatim}

with actual hash from:

\begin{verbatim}
slappasswd
\end{verbatim}

Add user:

\begin{verbatim}
ldapadd -x \
-D "cn=Manager,dc=hpc,dc=local" \
-W \
-f ~/user.ldif
\end{verbatim}
\subsection{Important LDAP User Attributes}
\label{sec:orgc387b9d}

\begin{center}
\begin{tabular}{ll}
Attribute & Purpose\\
\hline
uid & username\\
uidNumber & UNIX UID\\
gidNumber & primary group\\
homeDirectory & user home\\
loginShell & login shell\\
cn & common name\\
sn & surname\\
\end{tabular}
\end{center}
\subsection{Verify LDAP User}
\label{sec:orge729288}

Run ONLY on:

\begin{verbatim}
master
\end{verbatim}

Verify user:

\begin{verbatim}
ldapsearch -x \
-b "dc=hpc,dc=local" uid=testUser
\end{verbatim}

Expected:
\begin{itemize}
\item LDAP user entry
\end{itemize}
\subsection{Configure LDAP Client on Login Node}
\label{sec:org44bde15}

Now begin client-side integration.

Run ONLY on:

\begin{verbatim}
login
\end{verbatim}

Install LDAP client packages:

\begin{verbatim}
sudo dnf install -y \
openldap-clients \
nss-pam-ldapd
\end{verbatim}
\subsection{Configure nslcd}
\label{sec:orgbce268d}

Run ONLY on:

\begin{verbatim}
login
\end{verbatim}

Create configuration:

\begin{verbatim}
sudo vim /etc/nslcd.conf
\end{verbatim}

Paste:

\begin{verbatim}
uid nslcd
gid ldap

uri ldap://master
base dc=hpc,dc=local
\end{verbatim}

Explanation:

\begin{center}
\begin{tabular}{ll}
Parameter & Purpose\\
\hline
uri & LDAP server\\
base & LDAP search root\\
\end{tabular}
\end{center}

Secure configuration:

\begin{verbatim}
sudo chmod 600 /etc/nslcd.conf
\end{verbatim}
\subsection{Enable nslcd}
\label{sec:org41398ae}

Run ONLY on:

\begin{verbatim}
login
\end{verbatim}

Enable service:

\begin{verbatim}
sudo systemctl enable --now nslcd
\end{verbatim}

Verify:

\begin{verbatim}
systemctl status nslcd
\end{verbatim}

Expected:

\begin{verbatim}
active (running)
\end{verbatim}
\subsection{Verify LDAP Connectivity}
\label{sec:org81c99fb}

Run ONLY on:

\begin{verbatim}
login
\end{verbatim}

Test LDAP communication:

\begin{verbatim}
ldapsearch -x \
-H ldap://master \
-b "dc=hpc,dc=local" uid=testUser
\end{verbatim}

Expected:
\begin{itemize}
\item LDAP user entry
\end{itemize}

This verifies:
\begin{itemize}
\item login node can communicate with master LDAP server
\end{itemize}
\subsection{Configure NSS}
\label{sec:org8b8e346}

Run ONLY on:

\begin{verbatim}
login
\end{verbatim}

Edit:

\begin{verbatim}
sudo vim /etc/nsswitch.conf
\end{verbatim}

Modify lines:

\begin{verbatim}
passwd: files ldap
group: files ldap
shadow: files ldap
\end{verbatim}

Important:

\begin{verbatim}
This tells Linux:
check local files first, then LDAP.
\end{verbatim}
\subsection{Restart nslcd}
\label{sec:org83334fe}

Run ONLY on:

\begin{verbatim}
login
\end{verbatim}

Restart service:

\begin{verbatim}
sudo systemctl restart nslcd
\end{verbatim}
\subsection{Verify LDAP User Resolution}
\label{sec:orgb890c8a}

Run ONLY on:

\begin{verbatim}
login
\end{verbatim}

Test centralized user lookup:

\begin{verbatim}
getent passwd testUser
\end{verbatim}

Expected:

\begin{verbatim}
testUser:x:10001:10000:Test User:/home/testUser:/bin/bash
\end{verbatim}

Important concept:

\begin{verbatim}
This proves Linux itself recognizes LDAP users.
\end{verbatim}

This is one of the most important LDAP verification steps.
\subsection{NOTE (Do this for master too)}
\label{sec:org066bc95}
Now we must configure: \textbf{master} with EXACT same setup (master should be able to recognise the user).
\subsection{Install Automatic Home Directory Support}
\label{sec:orgd3a9ff2}

Run ONLY on:

\begin{verbatim}
login
\end{verbatim}

Install package:

\begin{verbatim}
sudo dnf install -y oddjob-mkhomedir
\end{verbatim}

Enable service:

\begin{verbatim}
sudo systemctl enable --now oddjobd
\end{verbatim}

Purpose:

\begin{verbatim}
Automatically creates user home directory
during first successful login.
\end{verbatim}
\subsection{Configure PAM Authentication}
\label{sec:org47b3e62}

Run ONLY on:

\begin{verbatim}
login
\end{verbatim}

Edit:

\begin{verbatim}
sudo vim /etc/pam.d/system-auth
\end{verbatim}

Add LDAP lines in appropriate sections.

Add in auth section:

\begin{verbatim}
auth        sufficient    pam_ldap.so use_first_pass
\end{verbatim}

Add in account section:

\begin{verbatim}
account     sufficient    pam_ldap.so
\end{verbatim}

Add in password section:

\begin{verbatim}
password    sufficient    pam_ldap.so use_authtok
\end{verbatim}

Add in session section:

\begin{verbatim}
session     optional      pam_ldap.so
session     optional      pam_mkhomedir.so
\end{verbatim}

Important:

\begin{verbatim}
Do NOT remove pam_unix.so lines.
\end{verbatim}
\subsection{Configure Second PAM File}
\label{sec:orgebad77f}

Run ONLY on:

\begin{verbatim}
login
\end{verbatim}

Edit:

\begin{verbatim}
sudo vim /etc/pam.d/password-auth
\end{verbatim}

Add SAME LDAP PAM lines.

Reason:

\begin{verbatim}
Different login methods use different PAM stacks.
\end{verbatim}
\subsection{Restart Services}
\label{sec:org2d17195}

Run ONLY on:

\begin{verbatim}
login
\end{verbatim}

Restart services:

\begin{verbatim}
sudo systemctl restart nslcd
sudo systemctl restart oddjobd
\end{verbatim}
\subsection{Verify LDAP Authentication}
\label{sec:orge150ca3}

Run ONLY on:

\begin{verbatim}
login
\end{verbatim}

Test login:

\begin{verbatim}
su - testUser
\end{verbatim}

Enter LDAP password.

Expected:

\begin{verbatim}
Creating directory '/home/testUser'
\end{verbatim}

Successful login indicates:

\begin{itemize}
\item LDAP authentication working
\item PAM integration working
\item automatic home directory creation working
\end{itemize}
\subsection{Verify User Home}
\label{sec:org808e074}

Run ONLY on:

\begin{verbatim}
login
\end{verbatim}

After login:

\begin{verbatim}
pwd
\end{verbatim}

Expected:

\begin{verbatim}
/home/testUser
\end{verbatim}
\begin{enumerate}
\item NOW WE MUST REPLICATE TO COMPUTE NODES
\label{sec:org423d84e}

Currently:

\begin{verbatim}
login node works fully
\end{verbatim}

Now we configure:

\begin{itemize}
\item compute01
\item compute02
\item compute03
\end{itemize}

with EXACT same client setup.
\end{enumerate}
\subsection{Troubleshooting}
\label{sec:org03ad2b9}
\begin{enumerate}
\item Important Slurm + LDAP Troubleshooting
\label{sec:orgbfe9147}

This section explains one of the MOST common HPC setup issues.
\item Symptom
\label{sec:org7c46844}

LDAP users:

\begin{itemize}
\item can SSH login
\item can use su
\item appear in getent passwd
\end{itemize}

BUT:

\begin{verbatim}
sbatch submit.sh
\end{verbatim}

results in:

\begin{itemize}
\item jobs stuck in pending
\item BeginTime state
\item scheduler instability
\item sinfo hanging
\item Slurm credential errors
\item slurmctld crashes
\end{itemize}
\item Root Cause
\label{sec:org7b24692}

Slurm internally works using:

\begin{itemize}
\item UID
\item GID
\item NSS identity resolution
\end{itemize}

When launching jobs, Slurm performs internal user identity lookups.

If master node cannot resolve LDAP users correctly through:

\begin{itemize}
\item NSS
\item nslcd
\item LDAP
\end{itemize}

then:

\begin{itemize}
\item Slurm credential generation fails
\item scheduler may crash
\item jobs cannot launch
\end{itemize}

Common error:

\begin{verbatim}
_fill_cred_gids: getpwuid failed for uid=10001
\end{verbatim}

Meaning:

\begin{verbatim}
Slurm cannot resolve LDAP user identity correctly.
\end{verbatim}
\item Important Verification Commands
\label{sec:orga382553}

Run on:

\begin{verbatim}
master
login
compute nodes
\end{verbatim}

Verify LDAP resolution:

\begin{verbatim}
getent passwd testUser
id testUser
\end{verbatim}

Expected:

\begin{verbatim}
uid=10001(testUser) gid=10000(hpcusers)
\end{verbatim}

These commands MUST work correctly on ALL cluster nodes.
\item Important Master Node Verification
\label{sec:org4a4f589}

The master node MUST also resolve LDAP users correctly.

Verify:

\begin{verbatim}
getent passwd testUser
\end{verbatim}

If master cannot resolve LDAP users:

\begin{itemize}
\item slurmctld may fail
\item scheduler may crash
\item jobs remain pending
\end{itemize}

Verify nslcd configuration on master:

\begin{verbatim}
sudo vim /etc/nslcd.conf
\end{verbatim}

Example:

\begin{verbatim}
uid nslcd
gid ldap

uri ldap://master
base dc=hpc,dc=local
\end{verbatim}

Verify NSS configuration:

\begin{verbatim}
sudo vim /etc/nsswitch.conf
\end{verbatim}

Ensure:

\begin{verbatim}
passwd: files ldap
group: files ldap
shadow: files ldap
\end{verbatim}

Restart LDAP client:

\begin{verbatim}
sudo systemctl restart nslcd
\end{verbatim}
\item Important Slurm Architecture
\label{sec:orge847df0}

Correct architecture:

\begin{center}
\begin{tabular}{ll}
Node & Service\\
\hline
master & slurmctld\\
login & login services only\\
compute01 & slurmd\\
compute02 & slurmd\\
compute03 & slurmd\\
\end{tabular}
\end{center}

IMPORTANT:

\begin{verbatim}
Login node should NOT run slurmd.
\end{verbatim}

Because login node is NOT configured as a compute node in:

\begin{verbatim}
/etc/slurm/slurm.conf
\end{verbatim}

Attempting to start slurmd on login gives:

\begin{verbatim}
slurmd: fatal: Unable to determine this slurmd's NodeName
\end{verbatim}

This is NORMAL and expected.

Correct action on login node:

\begin{verbatim}
sudo systemctl stop slurmd
sudo systemctl disable slurmd
\end{verbatim}

Expected afterward:

\begin{verbatim}
inactive (dead)
\end{verbatim}

This is CORRECT for login nodes.
\item Important Slurm Recovery Commands
\label{sec:org4b00551}

Whenever:

\begin{itemize}
\item LDAP users are added
\item NSS configuration changes
\item PAM configuration changes
\item nslcd configuration changes
\end{itemize}

restart Slurm services.

Restart controller on master:

\begin{verbatim}
sudo systemctl restart slurmctld
\end{verbatim}

Restart compute daemons on compute nodes:

\begin{verbatim}
sudo systemctl restart slurmd
\end{verbatim}

Restart LDAP client:

\begin{verbatim}
sudo systemctl restart nslcd
\end{verbatim}

Verify cluster:

\begin{verbatim}
sinfo
\end{verbatim}

Expected:

\begin{verbatim}
debug* up infinite 3 idle compute[01-03]
\end{verbatim}
\item Important Scheduler Debugging
\label{sec:org5d409aa}

If:

\begin{verbatim}
sinfo
\end{verbatim}

hangs indefinitely:

Possible cause:

\begin{verbatim}
slurmctld crashed
\end{verbatim}

Verify:

\begin{verbatim}
systemctl status slurmctld
\end{verbatim}

or:

\begin{verbatim}
journalctl -u slurmctld -n 50
\end{verbatim}
\end{enumerate}
\subsection{Current Cluster State After LDAP Setup}
\label{sec:orgb1a3089}

Cluster now contains:

\begin{center}
\begin{tabular}{ll}
Component & Status\\
\hline
OpenLDAP server & working\\
centralized users & working\\
centralized authentication & working\\
nslcd integration & working\\
NSS integration & working\\
PAM authentication & working\\
automatic home creation & working\\
\end{tabular}
\end{center}
\subsection{Important Debugging Commands}
\label{sec:org991ef45}

Verify LDAP tree:

\begin{verbatim}
ldapsearch -x -b "dc=hpc,dc=local"
\end{verbatim}

Verify LDAP user:

\begin{verbatim}
ldapsearch -x -b "dc=hpc,dc=local" uid=testUser
\end{verbatim}

Verify LDAP authentication:

\begin{verbatim}
ldapwhoami -x \
-D "uid=testUser,ou=People,dc=hpc,dc=local" \
-W \
-H ldap://master
\end{verbatim}

Verify Linux user resolution:

\begin{verbatim}
getent passwd testUser
\end{verbatim}

Verify nslcd service:

\begin{verbatim}
systemctl status nslcd
\end{verbatim}
\subsection{Important Architecture Understanding}
\label{sec:orgab04548}

Master node:
\begin{itemize}
\item provides centralized identity database
\end{itemize}

Login/compute nodes:
\begin{itemize}
\item authenticate users
\item resolve LDAP users
\item create user sessions
\end{itemize}

This separation is fundamental in distributed authentication systems.
\section{Automated LDAP User Creation Script}
\label{sec:org2a907f3}

Managing users manually using individual LDIF files becomes difficult
as the number of users increases.

To simplify user management, an automation script can be used to:

\begin{itemize}
\item create LDAP group
\item create LDAP user
\item generate password automatically
\item generate password hash automatically
\item add entries into LDAP
\end{itemize}

This script automates complete LDAP user creation.

IMPORTANT:

\begin{verbatim}
The script creates:

- private group for user
- LDAP user entry
- hashed LDAP password

automatically.
\end{verbatim}
\subsection{Script Location}
\label{sec:orgd0488d5}

Example:

\begin{verbatim}
cd ~/ldap-scripts
ls
\end{verbatim}

Expected:

\begin{verbatim}
createUser.sh
\end{verbatim}
\subsection{User Creation Script}
\label{sec:org2327f3d}

Create file:

\begin{verbatim}
vim createUser.sh
\end{verbatim}

Paste the complete script exactly as shown below:

\begin{verbatim}
#!/bin/bash 

LDAP_ADMIN="cn=Manager,dc=hpc,dc=local"
BASE_DN="dc=hpc,dc=local"

USERNAME=$1
UIDNUM=$2
GIDNUM=$3

if [ $# -ne 3 ]; then
    echo "Usage: $0 <username> <uid> <gid>"
    exit 1
fi

# Check if user already exists
USER_EXISTS=$(ldapsearch -x \
-b "ou=People,${BASE_DN}" "(uid=${USERNAME})" dn \
| grep "^dn:")

if [ ! -z "$USER_EXISTS" ]; then
    echo "ERROR: User ${USERNAME} already exists!"
    exit 1
fi

# Check if group already exists
GROUP_EXISTS=$(ldapsearch -x \
-b "ou=Groups,${BASE_DN}" "(cn=${USERNAME})" dn \
| grep "^dn:")

if [ ! -z "$GROUP_EXISTS" ]; then
    echo "ERROR: Group ${USERNAME} already exists!"
    exit 1
fi

PASSWORD="${USERNAME}@@123"

PASSWORD_HASH=$(slappasswd -s "$PASSWORD")

echo
echo "=================================="
echo "Creating LDAP User"
echo "=================================="
echo "Username : ${USERNAME}"
echo "UID      : ${UIDNUM}"
echo "GID      : ${GIDNUM}"
echo "=================================="

# Create private group LDIF
cat > /tmp/${USERNAME}_group.ldif <<EOF
dn: cn=${USERNAME},ou=Groups,${BASE_DN}
objectClass: posixGroup
cn: ${USERNAME}
gidNumber: ${GIDNUM}
EOF

# Create user LDIF
cat > /tmp/${USERNAME}.ldif <<EOF
dn: uid=${USERNAME},ou=People,${BASE_DN}
objectClass: inetOrgPerson
objectClass: posixAccount
objectClass: shadowAccount
cn: ${USERNAME}
sn: ${USERNAME}
uid: ${USERNAME}
uidNumber: ${UIDNUM}
gidNumber: ${GIDNUM}
homeDirectory: /home/${USERNAME}
loginShell: /bin/bash
userPassword: ${PASSWORD_HASH}
EOF

echo
echo "Adding LDAP group..."

ldapadd -x \
-D "${LDAP_ADMIN}" \
-W \
-f /tmp/${USERNAME}_group.ldif

if [ $? -ne 0 ]; then
    echo "ERROR: Failed to create LDAP group!"
    exit 1
fi

echo
echo "Adding LDAP user..."

ldapadd -x \
-D "${LDAP_ADMIN}" \
-W \
-f /tmp/${USERNAME}.ldif

if [ $? -ne 0 ]; then
    echo "ERROR: Failed to create LDAP user!"
    exit 1
fi

echo
echo "=================================="
echo "LDAP User Created Successfully"
echo "=================================="
echo "Username : ${USERNAME}"
echo "Password : ${PASSWORD}"
echo "UID      : ${UIDNUM}"
echo "GID      : ${GIDNUM}"
echo "=================================="
\end{verbatim}
\subsection{Make Script Executable}
\label{sec:orgbe65a99}

Before using the script:

\begin{verbatim}
chmod +x createUser.sh
\end{verbatim}
\subsection{Understanding Script Arguments}
\label{sec:org7518fd7}

The script requires THREE arguments:

\begin{center}
\begin{tabular}{ll}
Argument & Meaning\\
\hline
username & LDAP username\\
uid & UNIX UID\\
gid & UNIX GID\\
\end{tabular}
\end{center}

Syntax:

\begin{verbatim}
./createUser.sh <username> <uid> <gid>
\end{verbatim}

Example:

\begin{verbatim}
./createUser.sh student01 11001 11001
\end{verbatim}

Explanation:

\begin{center}
\begin{tabular}{rl}
Value & Meaning\\
\hline
student01 & username\\
11001 & UID\\
11001 & GID\\
\end{tabular}
\end{center}

IMPORTANT:

\begin{verbatim}
The script creates a private group automatically.

Therefore username and group name become identical.
\end{verbatim}

Example:

\begin{center}
\begin{tabular}{ll}
Entity & Created Automatically\\
\hline
User & student01\\
Group & student01\\
Home Directory & /home/student01\\
\end{tabular}
\end{center}
\subsection{Default Password Format}
\label{sec:org8acd33c}

The script automatically creates password using:

\begin{verbatim}
username@@123
\end{verbatim}

Example:

\begin{center}
\begin{tabular}{ll}
Username & Password\\
\hline
student01 & student01@@123\\
student02 & student02@@123\\
\end{tabular}
\end{center}

IMPORTANT:

\begin{verbatim}
Users should change password after first login.
\end{verbatim}
\subsection{What the Script Automatically Does}
\label{sec:org0dd06b3}

The script performs following operations automatically:

\begin{center}
\begin{tabular}{ll}
Operation & Description\\
\hline
Check existing user & avoids duplicate users\\
Check existing group & avoids duplicate groups\\
Generate password hash & secure LDAP password\\
Create LDAP group LDIF & automatic\\
Create LDAP user LDIF & automatic\\
Add LDAP group & automatic\\
Add LDAP user & automatic\\
\end{tabular}
\end{center}
\subsection{Example User Creation}
\label{sec:org470bc9d}

Example:

\begin{verbatim}
./createUser.sh student02 11002 11002
\end{verbatim}

Expected output:

\begin{verbatim}
==================================
Creating LDAP User
==================================
Username : student02
UID      : 11002
GID      : 11002
==================================
\end{verbatim}

LDAP admin password will be requested.

After successful creation:

\begin{verbatim}
==================================
LDAP User Created Successfully
==================================
Username : student02
Password : student02@@123
UID      : 11002
GID      : 11002
==================================
\end{verbatim}
\subsection{Verify User Creation}
\label{sec:orgd7f785e}

Verify user using:

\begin{verbatim}
getent passwd student02
\end{verbatim}

or:

\begin{verbatim}
id student02
\end{verbatim}

Expected:

\begin{verbatim}
uid=11002(student02)
\end{verbatim}

Verify LDAP entry:

\begin{verbatim}
ldapsearch -x \
-b "dc=hpc,dc=local" uid=student02
\end{verbatim}
\subsection{Verify Login}
\label{sec:org3469cae}

Test login:

\begin{verbatim}
su - student02
\end{verbatim}

or:

\begin{verbatim}
ssh student02@login
\end{verbatim}

Expected:

\begin{verbatim}
successful login
\end{verbatim}
\subsection{Important UID/GID Notes}
\label{sec:orgca1eb36}

IMPORTANT:

\begin{verbatim}
UID and GID values must remain unique cluster-wide.
\end{verbatim}

Avoid:
\begin{itemize}
\item duplicate UID
\item duplicate GID
\end{itemize}

because duplicate IDs can create:

\begin{itemize}
\item ownership conflicts
\item wrong usernames
\item incorrect file ownership
\item Slurm identity problems
\end{itemize}

Example of BAD configuration:

\begin{center}
\begin{tabular}{lr}
User & UID\\
\hline
student01 & 11001\\
student02 & 11001\\
\end{tabular}
\end{center}

This can cause:

\begin{verbatim}
whoami returning wrong username
\end{verbatim}

Always use unique UID/GID values.
\subsection{Important HPC Identity Concept}
\label{sec:orgc48b622}

In HPC clusters:

\begin{itemize}
\item NFS
\item Slurm
\item MPI
\item LDAP
\item SSH
\end{itemize}

all depend on:

\begin{verbatim}
consistent unique UID/GID mappings
\end{verbatim}

across the entire cluster.

This is one of the MOST important concepts in centralized HPC authentication.
\section{Automated LDAP User Deletion Script}
\label{sec:org7e472cc}

Managing users manually using:

\begin{itemize}
\item ldapdelete
\item group deletion
\item home directory cleanup
\end{itemize}

can become repetitive and error-prone.

To simplify cleanup operations, an automated LDAP deletion script can be used.

This script automates:

\begin{itemize}
\item LDAP user deletion
\item optional LDAP group deletion
\item optional home directory deletion
\end{itemize}

IMPORTANT:

\begin{verbatim}
The script provides interactive cleanup options
for safer user management.
\end{verbatim}
\subsection{Script Location}
\label{sec:org143ac42}

Example:

\begin{verbatim}
cd ~/ldap-scripts
ls
\end{verbatim}

Expected:

\begin{verbatim}
deleteUser.sh
\end{verbatim}
\subsection{User Deletion Script}
\label{sec:orgbf78e08}

Create file:

\begin{verbatim}
vim deleteUser.sh
\end{verbatim}

Paste the complete script exactly as shown below:

\begin{verbatim}
#!/bin/bash 

LDAP_ADMIN="cn=Manager,dc=hpc,dc=local"
BASE_DN="dc=hpc,dc=local"

USERNAME=$1

if [ $# -ne 1 ]; then
    echo "Usage: $0 <username>"
    exit 1
fi

echo
echo "=================================="
echo "Deleting LDAP User"
echo "=================================="
echo "Username : ${USERNAME}"
echo "=================================="

echo
echo "Deleting LDAP user entry..."

ldapdelete -x \
-D "${LDAP_ADMIN}" \
-W \
"uid=${USERNAME},ou=People,${BASE_DN}"

if [ $? -ne 0 ]; then
    echo "ERROR: Failed to delete LDAP user!"
    exit 1
fi

echo
echo "LDAP user deleted successfully."

# Ask about deleting group
echo
read -p "Delete LDAP group '${USERNAME}'? (y/n): " DELETE_GROUP

case $DELETE_GROUP in
    y|Y)

        ldapdelete -x \
        -D "${LDAP_ADMIN}" \
        -W \
        "cn=${USERNAME},ou=Groups,${BASE_DN}"

        if [ $? -eq 0 ]; then
            echo "LDAP group deleted successfully."
        else
            echo "WARNING: Failed to delete LDAP group."
        fi
        ;;

    n|N)
        echo "Skipping LDAP group deletion."
        ;;

    *)
        echo "Invalid option. Skipping group deletion."
        ;;
esac

# Ask about deleting home directory
echo
read -p "Delete home directory '/home/${USERNAME}'? (y/n): " DELETE_HOME

case $DELETE_HOME in
    y|Y)

        if [ -d "/home/${USERNAME}" ]; then
            sudo rm -rf "/home/${USERNAME}"

            if [ $? -eq 0 ]; then
                echo "Home directory deleted successfully."
            else
                echo "WARNING: Failed to delete home directory."
            fi
        else
            echo "Home directory does not exist."
        fi
        ;;

    n|N)
        echo "Skipping home directory deletion."
        ;;

    *)
        echo "Invalid option. Skipping home deletion."
        ;;
esac

echo
echo "=================================="
echo "User cleanup completed"
echo "=================================="
\end{verbatim}
\subsection{Make Script Executable}
\label{sec:org21e0caa}

Before using the script:

\begin{verbatim}
chmod +x deleteUser.sh
\end{verbatim}
\subsection{Understanding Script Arguments}
\label{sec:orgcdf6a52}

The script requires ONE argument:

\begin{center}
\begin{tabular}{ll}
Argument & Meaning\\
\hline
username & LDAP username\\
\end{tabular}
\end{center}

Syntax:

\begin{verbatim}
./deleteUser.sh <username>
\end{verbatim}

Example:

\begin{verbatim}
./deleteUser.sh student02
\end{verbatim}
\subsection{What the Script Automatically Does}
\label{sec:orgb5df399}

The script performs following operations automatically:

\begin{center}
\begin{tabular}{ll}
Operation & Description\\
\hline
Delete LDAP user & removes LDAP account\\
Optional group deletion & removes LDAP group if requested\\
Optional home deletion & removes user home if requested\\
Interactive confirmation & prevents accidental cleanup\\
\end{tabular}
\end{center}
\subsection{Example User Deletion}
\label{sec:orgc3b3462}

Example:

\begin{verbatim}
./deleteUser.sh student02
\end{verbatim}

Expected output:

\begin{verbatim}
==================================
Deleting LDAP User
==================================
Username : student02
==================================
\end{verbatim}

LDAP admin password will be requested.

After successful deletion:

\begin{verbatim}
LDAP user deleted successfully.
\end{verbatim}

The script then asks:

\begin{verbatim}
Delete LDAP group 'student02'? (y/n):
\end{verbatim}

If:

\begin{verbatim}
y
\end{verbatim}

is selected, corresponding LDAP group will also be removed.

Next prompt:

\begin{verbatim}
Delete home directory '/home/student02'? (y/n):
\end{verbatim}

If:

\begin{verbatim}
y
\end{verbatim}

is selected, user home directory will also be deleted.
\subsection{Important Group Deletion Notes}
\label{sec:org0457d6f}

IMPORTANT:

\begin{verbatim}
Do NOT delete shared groups accidentally.
\end{verbatim}

Example:

\begin{verbatim}
hpcusers
\end{verbatim}

may contain multiple users.

Deleting shared groups can create:

\begin{itemize}
\item broken group mappings
\item permission problems
\item Slurm identity issues
\end{itemize}

The script therefore asks confirmation before deleting groups.
\subsection{Important Home Directory Notes}
\label{sec:org2fa46fa}

Deleting home directories permanently removes:

\begin{itemize}
\item source codes
\item job scripts
\item outputs
\item datasets
\item SSH keys
\item configuration files
\end{itemize}

IMPORTANT:

\begin{verbatim}
Home directory deletion is irreversible.
\end{verbatim}

Always verify before deleting.
\subsection{Verify User Removal}
\label{sec:org85a3d5e}

Verify deletion:

\begin{verbatim}
getent passwd student02
\end{verbatim}

Expected:

\begin{verbatim}
no output
\end{verbatim}

Verify LDAP entry removal:

\begin{verbatim}
ldapsearch -x \
-b "dc=hpc,dc=local" uid=student02
\end{verbatim}

Expected:

\begin{verbatim}
no entries returned
\end{verbatim}

Verify login failure:

\begin{verbatim}
su - student02
\end{verbatim}

Expected:

\begin{verbatim}
user does not exist
\end{verbatim}
\subsection{Important HPC User Management Concept}
\label{sec:orgd0ce4ac}

In centralized HPC environments:

\begin{itemize}
\item user lifecycle management
\item account cleanup
\item UID/GID consistency
\item group management
\end{itemize}

are critical administrative tasks.

Improper deletion can create:

\begin{itemize}
\item orphaned files
\item broken permissions
\item Slurm scheduling issues
\item incorrect ownership mappings
\end{itemize}

Therefore automated and controlled cleanup procedures are highly important in production HPC systems.
\section{Automated LDAP Password Change Script}
\label{sec:orgc1b764f}

Managing passwords manually using LDIF files becomes inconvenient
in centralized LDAP environments.

To simplify password management, an automated password change script
can be used.

This script automates:

\begin{itemize}
\item LDAP password update
\item password hash generation
\item secure password input
\item LDAP password modification
\end{itemize}

IMPORTANT:

\begin{verbatim}
Because authentication is centralized through LDAP,
password changes automatically work cluster-wide.
\end{verbatim}

After changing password:

\begin{itemize}
\item login node recognizes new password
\item compute nodes recognize new password
\item master recognizes new password
\end{itemize}

without changing anything individually on nodes.
\subsection{Script Location}
\label{sec:org692b9c8}

Example:

\begin{verbatim}
cd ~/ldap-scripts
ls
\end{verbatim}

Expected:

\begin{verbatim}
changePassword.sh
\end{verbatim}
\subsection{Password Change Script}
\label{sec:org0e1f009}

Create file:

\begin{verbatim}
vim changePassword.sh
\end{verbatim}

Paste the complete script exactly as shown below:

\begin{verbatim}
#!/bin/bash

LDAP_ADMIN="cn=Manager,dc=hpc,dc=local"
BASE_DN="dc=hpc,dc=local"

USERNAME=$1

if [ $# -ne 1 ]; then
    echo "Usage: $0 <username>"
    exit 1
fi

# Check if user exists
USER_EXISTS=$(ldapsearch -x \
-b "ou=People,${BASE_DN}" "(uid=${USERNAME})" dn \
| grep "^dn:")

if [ -z "$USER_EXISTS" ]; then
    echo "ERROR: User ${USERNAME} does not exist!"
    exit 1
fi

echo
echo "=================================="
echo "LDAP Password Change"
echo "=================================="
echo "Username : ${USERNAME}"
echo "=================================="

# Take password securely
echo
read -s -p "Enter new password: " PASSWORD1
echo

read -s -p "Confirm new password: " PASSWORD2
echo

# Check password match
if [ "$PASSWORD1" != "$PASSWORD2" ]; then
    echo
    echo "ERROR: Passwords do not match!"
    exit 1
fi

# Generate LDAP password hash
PASSWORD_HASH=$(slappasswd -s "$PASSWORD1")

# Create LDIF file
cat > /tmp/${USERNAME}_password.ldif <<EOF
dn: uid=${USERNAME},ou=People,${BASE_DN}
changetype: modify
replace: userPassword
userPassword: ${PASSWORD_HASH}
EOF

echo
echo "Updating LDAP password..."

ldapmodify -x \
-D "${LDAP_ADMIN}" \
-W \
-f /tmp/${USERNAME}_password.ldif

if [ $? -ne 0 ]; then
    echo
    echo "ERROR: Failed to change password!"
    exit 1
fi

echo
echo "=================================="
echo "Password updated successfully"
echo "=================================="
echo "Username : ${USERNAME}"
echo "=================================="
\end{verbatim}
\subsection{Make Script Executable}
\label{sec:org0d503dd}

Before using the script:

\begin{verbatim}
chmod +x changePassword.sh
\end{verbatim}
\subsection{Understanding Script Arguments}
\label{sec:orga639b05}

The script requires ONE argument:

\begin{center}
\begin{tabular}{ll}
Argument & Meaning\\
\hline
username & LDAP username\\
\end{tabular}
\end{center}

Syntax:

\begin{verbatim}
./changePassword.sh <username>
\end{verbatim}

Example:

\begin{verbatim}
./changePassword.sh student01
\end{verbatim}
\subsection{What the Script Automatically Does}
\label{sec:org18847b0}

The script performs following operations automatically:

\begin{center}
\begin{tabular}{ll}
Operation & Description\\
\hline
Verify LDAP user exists & prevents invalid modifications\\
Secure password input & password hidden during typing\\
Password confirmation & avoids typing mistakes\\
Generate password hash & secure LDAP password storage\\
Modify LDAP password & centralized password update\\
\end{tabular}
\end{center}
\subsection{Example Password Change}
\label{sec:orgd1252c9}

Example:

\begin{verbatim}
./changePassword.sh student01
\end{verbatim}

Expected output:

\begin{verbatim}
==================================
LDAP Password Change
==================================
Username : student01
==================================
\end{verbatim}

The script asks:

\begin{verbatim}
Enter new password:
\end{verbatim}

Then:

\begin{verbatim}
Confirm new password:
\end{verbatim}

LDAP admin password will also be requested.

After successful update:

\begin{verbatim}
==================================
Password updated successfully
==================================
Username : student01
==================================
\end{verbatim}
\subsection{Verify Password Change}
\label{sec:org2f492df}

Test login using updated password:

\begin{verbatim}
su - student01
\end{verbatim}

or:

\begin{verbatim}
ssh student01@login
\end{verbatim}

Expected:

\begin{verbatim}
successful login using new password
\end{verbatim}

Because LDAP authentication is centralized, the updated password works automatically on:

\begin{itemize}
\item login node
\item compute nodes
\item master node
\end{itemize}
\subsection{Important Password Management Notes}
\label{sec:orga0d0e16}

IMPORTANT:

\begin{verbatim}
LDAP stores password hashes, not plain-text passwords.
\end{verbatim}

The script automatically generates secure LDAP-compatible password hashes using:

\begin{verbatim}
slappasswd
\end{verbatim}

This improves security by avoiding direct plain-text password storage.
\end{document}
